\chapter{Walks, Chains and Groups}

\begin{center}\itshape
Chance is only the measure of our ignorance.\\
\upshape --- Henri Poincaré, ``Chance'', in \emph{Science and Method} (1908), English translation by Francis Maitland (1914)
\end{center}

In 1872 Ludwig Boltzmann proved his $H$-theorem. Take a gas, describe it by the distribution of the velocities of its molecules, and
follow that distribution as the molecules collide: a certain quantity $H$, built from the distribution, can only decrease, until the
gas reaches the equilibrium of Maxwell. The arrow of time, the fact that a drop of ink spreads in water and never gathers itself
again, seemed to come out of Newton's laws.

Two objections followed, and both are about walks and groups. In 1876 Josef Loschmidt, Boltzmann's friend and colleague in Vienna,
pointed out that Newton's laws do not care about the direction of time: reverse every velocity at some instant, and the gas
retraces its history, with $H$ going \emph{up}. In 1896 Ernst Zermelo, then an assistant of Max Planck in Berlin, used a theorem that
Henri Poincaré had proved in 1890: a mechanical system confined to a bounded region comes back, after a long enough time,
arbitrarily close to its initial state. So $H$ cannot decrease for ever; it must come back up.

Boltzmann's answer was that his theorem is a statement about what happens \emph{typically}, not always. Sixty years later Mark Kac
built a toy that makes this answer precise. Put $n$ balls, black or white, on the sites of a circle, and mark some of the $n$ edges
between neighbouring sites at random. At each tick every ball moves one site clockwise, and a ball that crosses a marked edge changes
colour. The rule is deterministic and can be run backwards; after $2n$ ticks every ball has crossed every edge twice, so the
initial state comes back exactly, as Poincaré and Zermelo demand. And yet, for a large ring, the proportion of black balls drifts
towards one half at an exponential rate and stays there for a very long time, with fluctuations that shrink as $n$ grows, as
Boltzmann demands. You will prove both halves of this reconciliation in Problem~\ref{pb:kacring}.

The ring is a walk on a finite group (the rotations of the circle) dressed with a little randomness, and that is the theme of the
chapter: walks and chains meeting the groups of Tome~2, the arithmetic of Tome~1 and the analysis of Tome~4. Groups decide when
a walk equidistributes, and symmetry computes hitting times; electrical networks turn hitting times into linear algebra; Markov
chains hide in the Collatz map and in the carries of school addition; and a few chains with a story, from three gamblers to the
genetics of a small population, close the chapter.

\begin{think}
A ring has $6$ sites, all carrying black balls, and $3$ of its $6$ edges are marked. Without computing, guess the colours after $6$
ticks, and after $12$. Now take a ring of $1000$ sites in which each edge is marked independently with probability $\frac1{10}$, all balls
black at the start. Roughly what proportion of the balls do you expect to be black after $10$ ticks? After $1000$ ticks?
\end{think}

Throughout the chapter, a \emph{random walk} on a finite group $G$ driven by a probability $\mu$ on $G$ is the Markov chain
$X_{t+1}=X_t\xi_{t+1}$, where $\xi_1,\xi_2,\dots$ are independent with law $\mu$; its transition matrix is $P(g,h)=\mu(g^{-1}h)$, and the
uniform law is stationary because every column of $P$ also sums to $1$. The walk is \emph{symmetric} if $\mu(g^{-1})=\mu(g)$. On a
graph, the simple random walk moves to a uniform neighbour; its stationary law is $\pi(x)=\frac{\deg x}{2|\mathcal E|}$ (Tome~3,
Corollary~8.15). We write $T_y$ for the first time the walk hits $y$ (at time $\ge1$ if it starts at $y$), and $\mathbb E_x$ for the
expectation when the walk starts at $x$.

\section{Walks on finite groups}

The first problems ask the questions of Tome~3, Chapter~8 (when does a walk equidistribute, and how fast) for walks on groups, where
the answers become algebra: cosets decide periodicity, characters compute the law, and the diameter bounds the speed.

\begin{exercise}[$\star\star$]\label{pb:returnbound}
\emph{The identity is the likeliest place.} Let $X$ be a symmetric random walk on a finite group $G$, started at the identity $e$,
and write $p_t(g)=P(X_t=g)$.
\begin{enumerate}[label=(\alph*)]
\item Show that $p_{2t}(e)=\sum_{g}p_t(g)^2$, and deduce that $p_{2t}(g)\le p_{2t}(e)$ for every $g\in G$ and that $p_{2t}(e)\ge\frac1{|G|}$.
\item Show that $t\mapsto p_{2t}(e)$ is non-increasing.
\item Show that $p_{2t+1}(e)$ can be $0$ for every $t$, but that $p_{2t+1}(g)\le p_{2t}(e)$ for all $g$ and $t$.
\item The walk on $\Z/n\Z$ with steps $\pm1$, started at $0$, is at time $2t$ at $0$ with probability $\frac1n\sum_{k=0}^{n-1}\cos^{2t}\frac{2\pi k}n$. Check (a) and
(b) directly on this formula.
\end{enumerate}
\end{exercise}

\begin{exercise}[$\star\star\star$]\label{pb:period}
\emph{Cosets decide periodicity.} Let $\mu$ be a probability on a finite group $G$ whose support $S$ generates $G$, and let $X$ be the
random walk driven by $\mu$, started at $e$.
\begin{enumerate}[label=(\alph*)]
\item Show that the walk is irreducible.
\item Suppose that $S\subseteq gN$ for a normal subgroup $N\ne G$ and some $g\in G$. Show that $G/N$ is cyclic, generated by $gN$, and that the
walk has period at least $[G:N]$; in particular it does not converge to the uniform law.
\item Conversely, let $d$ be the period of the walk and $N$ the set of elements that the walk can reach at times divisible by
$d$. Show that $N$ is a normal subgroup of index $d$ and that $S\subseteq sN$ for any $s\in S$. Conclude: \emph{the walk converges to the
uniform law if and only if $S$ is contained in no coset of a proper normal subgroup.}
\item A random shuffle applies a uniformly chosen transposition to a deck of $n\ge2$ cards. Show that it never converges, and
compute its period.
\item Let $p\ge5$ be prime, and multiply a residue $x\in(\Z/p\Z)^*$ at each step by $2$ or by $3$, with probability $\frac12$ each. Assuming
that $2$ and $3$ generate $(\Z/p\Z)^*$, show that the walk converges to the uniform law if and only if $3\cdot2^{-1}$ is a primitive root
modulo $p$. Check that it fails for $p=5$ and $p=13$, and find the period in these two cases.
\end{enumerate}
\end{exercise}

\begin{exercise}[$\star\star$]\label{pb:rollingdie}
\emph{A die rolling on a chessboard.} A cube sits on a square of an infinite chessboard, exactly covering it. A move tips the cube
over one of the four edges of its bottom face, onto one of the four neighbouring squares.
\begin{enumerate}[label=(\alph*)]
\item Each move rotates the cube by a quarter turn about a horizontal axis. Using the action of the rotations on the four long
diagonals of the cube, show that every move is an odd permutation of the diagonals.
\item The cube makes a sequence of moves and comes back to its starting square with the same face on top. Show that it cannot
have turned by a quarter turn about the vertical axis. Can it have turned by a half turn?
\item Show that rolling the cube around the four squares of a $2\times2$ block (east, north, west, south) rotates it by a third
of a turn about a long diagonal. Deduce that every orientation allowed by (b), that is every even permutation of the
diagonals, can be obtained by a closed tour that comes back to the starting square.
\item Now the cube makes random moves, each of the four directions with probability $\frac14$. Show that the orientation of the cube
does not converge in law, but that the face on top does: after $t$ moves, each face is on top with probability tending to $\frac16$.
\item What is the expected number of moves before the cube is again in its initial orientation (wherever it is on the board)?
\end{enumerate}
\end{exercise}

\begin{exercise}[$\star\star\star$]\label{pb:dodecaturns}
\emph{Random turns of an icosahedron.} The rotation group of the regular icosahedron is isomorphic to $A_5$ (Tome~2, Theorem~12.5), and the
$20$ rotations by a third of a turn about the axes through the centres of opposite faces form the conjugacy class of the
$3$-cycles. An icosahedron is turned, again and again, by a uniformly chosen one of these $20$ rotations.
\begin{enumerate}[label=(\alph*)]
\item Let $\mu$ be a probability on a finite group $G$ that is constant on conjugacy classes, and $\mu^{*t}$ the law of the walk at time
$t$, started at $e$. Show that, for every $g\in G$,
\[
\mu^{*t}(g)=\frac1{|G|}\sum_\chi\chi(1)\overline{\chi(g)}\,\Bigl(\frac{\hat\mu(\chi)}{\chi(1)}\Bigr)^t,\qquad\hat\mu(\chi)=\sum_{h\in G}\mu(h)\chi(h),
\]
the sum being over the irreducible characters of $G$. (Show first that, for a class function $f$ and an irreducible representation $\rho$ of
dimension $d$, the matrix $\sum_hf(h)\rho(h)$ is the scalar $\frac1d\sum_hf(h)\chi_\rho(h)$.)
\item Using the character table of $A_5$ (Tome~2, Exercise~11.7), show that the icosahedron is back in its initial position after $t$
turns with probability
\[
\frac1{60}\Bigl(1+16\cdot4^{-t}+25\cdot\Bigl(-\frac15\Bigr)^t\Bigr).
\]
\item Find the probability that after $t$ turns the icosahedron has turned by a fifth of a turn (in either direction) about some
axis through opposite vertices, and show that the law of the position converges to the uniform law at the rate $4^{-t}$.
\end{enumerate}
\end{exercise}

\begin{exercise}[$\star\star\star$]\label{pb:gapdiam}
\emph{The diameter controls the speed.} Let $G$ be a finite group, $S$ a generating set with $S=S^{-1}$, and $d$ the diameter of the
Cayley graph: every element is a product of at most $d$ elements of $S$. Let $P(g,gs)=\frac1{|S|}$ for $s\in S$ be the symmetric walk.
For $f:G\to\R$ let
\[
\mathcal E(f)=\frac1{2|G||S|}\sum_{g\in G}\sum_{s\in S}\bigl(f(gs)-f(g)\bigr)^2,\qquad\operatorname{Var}(f)=\frac1{2|G|^2}\sum_{g,h\in G}\bigl(f(h)-f(g)\bigr)^2.
\]
\begin{enumerate}[label=(\alph*)]
\item Show that $\mathcal E(f)=\langle f,(I-P)f\rangle$ and $\operatorname{Var}(f)=\langle f,f\rangle-\langle f,\mathbf 1\rangle^2$ for the inner product $\langle f,h\rangle=\frac1{|G|}\sum_gf(g)h(g)$, and deduce
that $1-\lambda_2=\min\frac{\mathcal E(f)}{\operatorname{Var}(f)}$ over non-constant $f$, where $\lambda_2$ is the second largest eigenvalue of $P$.
\item Show that $\operatorname{Var}(f)\le|S|\,d^2\,\mathcal E(f)$ for every $f$. (Fix, for each $x\in G$, a word of length at most $d$ for $x$, and write $f(gx)-f(g)$ as a
telescoping sum along it.) Conclude that $1-\lambda_2\ge\frac1{|S|d^2}$.
\item For $\Z/n\Z$ with $S=\{\pm1\}$, compare with the true value $1-\cos\frac{2\pi}n$. For $\mathfrak S_n$ generated by the $n-1$ adjacent transpositions,
what lower bound do you get?
\end{enumerate}
\end{exercise}

Figure~\ref{fig:treereturn} shows how differently two worlds with four roads at every crossing treat a walker who wants to come home.

\begin{exercise}[$\star\star\star$]\label{pb:treewalk}
\emph{A walk that escapes in a $4$-regular world.} The free group on $a,b$ is the set of reduced words in $a,a^{-1},b,b^{-1}$; its Cayley
graph is the infinite $4$-regular tree. A walk starts at the empty word $e$ and multiplies on the right by one of the four
letters, chosen uniformly. Let $D_t$ be the length of the reduced word at time $t$.
\begin{enumerate}[label=(\alph*)]
\item Show that $(D_t)$ is a Markov chain: from $0$ it goes to $1$, and from $k\ge1$ it goes to $k+1$ with probability $\frac34$ and to
$k-1$ with probability $\frac14$.
\item Show that the walk returns to $e$ with probability $\frac13$, and that the expected number of visits to $e$ (time $0$ included) is $\frac32$.
\item Let $h(k)=3^{-k/2}$. Show that $\mathbb E[h(D_{t+1})\mid D_t]\le\frac{\sqrt3}2h(D_t)$, and deduce that $P(D_{2n}=0)\le(\frac{\sqrt3}2)^{2n}$.
\item By counting the paths of $D$ that stay positive between times $1$ and $2n-1$, show that $P(D_{2n}=0)\ge\frac1{8n^2}(\frac34)^{n-1}$, and conclude that
$P(D_{2n}=0)^{1/2n}\to\frac{\sqrt3}2$.
\item Show that on the $q$-regular tree ($q\ge3$) the same limit is $\frac{2\sqrt{q-1}}q$. The simple walk on $\Z^2$ also has four neighbours at
every vertex: what is the limit there?
\end{enumerate}
\end{exercise}

\begin{figure}[ht]
\centering
\fig{ch11-treereturn}
\caption{The probability of being back at the start after $2k$ steps, on a logarithmic scale. On the $4$-regular tree of the free group
(Problem~\ref{pb:treewalk}) it decays like $(\frac34)^k$, up to a power of $k$; on $\Z^2$, where every point also has four neighbours, it decays only like
$\frac1{\pi k}$.}
\label{fig:treereturn}
\end{figure}

\section{Electricity and symmetry}

Tome~3, Exercise~7.8 turned the walk on a graph into an electrical network: with a unit resistor on each edge, the probability of
reaching $b$ before $a$ is a voltage, and the effective resistance $R(a,b)$ measures how hard it is to escape from $a$ to $b$. The
problems of this section push the dictionary further: resistances compute hitting times, symmetry computes resistances, and an
energy principle decides recurrence without a single Fourier integral. Throughout, $G$ is a finite connected graph with vertex set $V$,
$|V|=n$, and edge set $\mathcal E$, and $L=D-A$ is its Laplacian (degree matrix minus adjacency matrix).

\begin{exercise}[$\star\star\star$]\label{pb:commute}
\emph{The commute time identity.} Consider the simple random walk on $G$.
\begin{enumerate}[label=(\alph*)]
\item Let $a\in V$ and let $\tau\ge1$ be a stopping time with $X_\tau=a$ and $\mathbb E_a[\tau]<\infty$. Show that the expected number of visits to $x$ during
the times $0,1,\dots,\tau-1$ is $\pi(x)\,\mathbb E_a[\tau]$, for every $x\in V$.
\item Let $a\ne b$, and let $\tau$ be the first return to $a$ after the first visit to $b$. Using Tome~3, Exercise~7.8(b), show that the
walk visits $a$ on average $\deg(a)R(a,b)$ times before $\tau$, and deduce the commute time identity
\[
\mathbb E_a[T_b]+\mathbb E_b[T_a]=2|\mathcal E|\,R(a,b).
\]
\item Check the identity on a path of length $m$ between its two ends, and on a cycle of length $n$ between two neighbours.
\item Show that $\mathbb E_a[T_b]<2|\mathcal E|$ when $a$ and $b$ are neighbours, and deduce that the walk visits every vertex in expected time at most
$2|\mathcal E|(n-1)$, from any starting point. (Walk around a spanning tree.)
\end{enumerate}
\end{exercise}

\begin{exercise}[$\star\star\star$]\label{pb:foster}
\emph{Foster's theorem and symmetric graphs.}
\begin{enumerate}[label=(\alph*)]
\item Let $a\ne b$. Show that there is a vector $v\in\R^V$ with $Lv=e_a-e_b$, unique up to adding a constant, and that
$R(a,b)=v(a)-v(b)$. (Relate $v$ to the voltage $h$ of Tome~3, Exercise~7.8.)
\item Prove Foster's theorem: $\sum_{\{x,y\}\in\mathcal E}R(x,y)=n-1$. (Work on the hyperplane $W=\mathbf 1^\perp$, where $L$ is invertible, and take a trace.)
\item Suppose the automorphism group of $G$ acts transitively on the edges. Show that $R(x,y)=\frac{n-1}{|\mathcal E|}$ for every edge. If moreover, for every edge
$\{x,y\}$, some automorphism exchanges $x$ and $y$, show that $\mathbb E_x[T_y]=n-1$ for neighbours $x,y$.
\item Deduce the expected time for a walker on the vertices of a cube, of the Petersen graph, and of a dodecahedron to reach a
given neighbour of its starting vertex.
\end{enumerate}
\end{exercise}

\begin{exercise}[$\star\star$]\label{pb:antipode}
\emph{Reaching the opposite corner.} A walker moves at random along the edges of a polyhedron, starting at a vertex $a$, and we want
the expected time to reach the opposite vertex $b$.
\begin{enumerate}[label=(\alph*)]
\item Suppose that for every vertex $x$ at distance $k$ from $b$, the numbers $c_k$, $a_k$, $b_k$ of neighbours of $x$ at distance $k-1$, $k$, $k+1$ from
$b$ depend only on $k$. Show that the distance $\operatorname{dist}(X_t,b)$ is a Markov chain. If $n_k$ is the number of vertices at distance $k$ from $b$ and $r$ the common
degree, show that the expected time to go from distance $k$ to distance $k-1$ is $\frac r{n_kc_k}\sum_{i\ge k}n_i$.
\item Show that $\mathbb E_a[T_b]$ is $6$ for the octahedron, $10$ for the cube, $15$ for the icosahedron and $35$ for the dodecahedron.
\item For the cube $\{0,1\}^m$ (vertices at Hamming distance $1$ are neighbours), show that the expected time from $0\cdots0$ to $1\cdots1$ is
\[
\sum_{k=1}^m\binom{m-1}{k-1}^{-1}\sum_{j=0}^{k-1}\binom mj,
\]
and that it is equivalent to $2^m$.
\item With a $1$-ohm resistor on every edge, find the resistance between opposite vertices of the cube, the icosahedron and the
dodecahedron. (Use Problem~\ref{pb:commute}.)
\end{enumerate}
\end{exercise}

\begin{exercise}[Journey]\label{pb:rayleigh}
\emph{Every subgraph of the plane grid is recurrent.} Let $a\ne b$ be vertices of $G$. A \emph{unit flow} from $a$ to $b$ is a function $\theta$ on
ordered pairs of neighbours with $\theta(x,y)=-\theta(y,x)$, $\sum_y\theta(x,y)=0$ for $x\ne a,b$, and $\sum_y\theta(a,y)=1$. Its \emph{energy} is
$\mathcal W(\theta)=\sum_{\{x,y\}\in\mathcal E}\theta(x,y)^2$.
\begin{enumerate}[label=(\alph*)]
\item With $v$ as in Problem~\ref{pb:foster}(a), show that $i(x,y)=v(x)-v(y)$ is a unit flow of energy $R(a,b)$, and that every unit flow has energy at least
$R(a,b)$ (Thomson's principle).
\item Deduce Rayleigh's monotonicity: removing edges from $G$ (keeping $a$ and $b$ connected) cannot decrease $R(a,b)$.
\item Let $\Pi_1,\dots,\Pi_m$ be disjoint sets of edges, each of which meets every path from $a$ to $b$. Show the Nash-Williams inequality
$R(a,b)\ge\sum_k\frac1{|\Pi_k|}$.
\item In $\Z^2$ let $B_n=\{-n,\dots,n\}^2$, and let $G_n$ be the graph obtained from $B_n$ (with the grid edges inside it) by merging its boundary into a
single vertex $z_n$. Show that $R(0,z_n)\ge\frac18\ln(2n+1)$.
\item Let $H$ be any set of edges of $\Z^2$ with at least one edge at $0$, and consider the walk that moves along the edges of $H$, to a
uniform $H$-neighbour. Show that it returns to $0$ with probability $1$.
\end{enumerate}
\end{exercise}

\begin{exercise}[Journey]\label{pb:aldousbroder}
\emph{A random walk draws a uniform spanning tree.} Run the simple random walk on $G$ from a vertex $r$ until it has visited every vertex,
and for each $v\ne r$ keep the edge along which the walk entered $v$ for the first time.
\begin{enumerate}[label=(\alph*)]
\item Show that the kept edges form a spanning tree, the \emph{first-entrance tree}.
\item An \emph{arborescence} rooted at $x$ is a spanning tree with every edge oriented towards $x$; each $u\ne x$ has one outgoing arc. The \emph{tree chain}
moves from $(T,x)$ as follows: pick a uniform neighbour $y$ of $x$, add the arc $x\to y$ and delete the arc leaving $y$. Show that the result is an
arborescence rooted at $y$. Show that if the tree chain follows the steps of a walk path that visits every vertex, its final state is the
\emph{last-exit tree} of the path, in which each vertex other than the endpoint points to the vertex visited just after its last visit.
\item Let $w(T,x)=\prod_{u\ne x}\frac1{\deg u}$. Show that $\mu(T,x)=\frac{w(T,x)}Z$ is a stationary law of the tree chain, that $\mu(T,x)=\frac{\deg x}{2|\mathcal E|\,\tau(G)}$ where $\tau(G)$
is the number of spanning trees, and that the second coordinate of a $\mu$-distributed state has law $\pi$.
\item Start the tree chain from $\mu$, run it for $N$ steps, and read the walk backwards (Tome~3, Proposition~8.14). Letting $N\to\infty$, prove the
theorem of Aldous and Broder: \emph{the first-entrance tree is uniformly distributed among the $\tau(G)$ spanning trees, whatever the starting
vertex $r$.}
\end{enumerate}
\end{exercise}

\section{Chains inside arithmetic}

Some Markov chains are planted in arithmetic by hand, like the walks modulo $p$ of earlier chapters. Others were there all along: in
the parities of the Collatz map, in the carries of school addition, in the coordinates of a lattice walk.

Let $T(m)=\frac m2$ if $m$ is even and $T(m)=\frac{3m+1}2$ if $m$ is odd. The Collatz conjecture says that every positive integer
eventually reaches $1$ under $T$; it is still open, and it comes back in Tome~15, Appendix~F. For $m\ge1$ and $k\ge1$ the
\emph{parity vector} of $m$ is $v_k(m)=(m\bmod2,\ T(m)\bmod2,\ \dots,\ T^{k-1}(m)\bmod2)\in\{0,1\}^k$.

\begin{exercise}[$\star\star$]\label{pb:terras}
\emph{Collatz parities are fair coins.}
\begin{enumerate}[label=(\alph*)]
\item Show that $v_k(m)$ depends only on $m$ modulo $2^k$, and that $m\bmod2^k\mapsto v_k(m)$ is a bijection from $\Z/2^k\Z$ onto $\{0,1\}^k$.
\item Deduce that if $m$ is uniform on $2^k$ consecutive integers, the parities of $m,T(m),\dots,T^{k-1}(m)$ are independent fair coin
flips.
\item Show that $T^k(m)=\lambda m+\rho$, where $\lambda=\frac{3^j}{2^k}$, $j$ is the number of ones in $v_k(m)$, and $\rho\ge0$ depends only on $v_k(m)$.
\item Find the smallest $m\ge1$ whose first five parities are all odd. The number $1$ has parity vector $(1,0,1,0,1)$ of length $5$; find the
next integer with the same vector.
\end{enumerate}
\end{exercise}

\begin{exercise}[$\star\star\star$]\label{pb:everett}
\emph{Almost every integer goes down.} Let $\sigma(m)=\min\{k\ge1:T^k(m)<m\}$, with $\sigma(m)=\infty$ if there is no such $k$.
\begin{enumerate}[label=(\alph*)]
\item Let $k\ge1$ and let $u\in\{0,1\}^k$ have $j$ ones with $3^j<2^k$. Show that $\sigma(m)\le k$ for all large enough $m$ with $v_k(m)=u$.
\item Deduce that the set of $m$ with $\sigma(m)\le k$ has lower density at least $P\bigl(B_k<k\frac{\ln2}{\ln3}\bigr)$, where $B_k$ is binomial with parameters $k$
and $\frac12$.
\item Conclude (Terras, 1976; Everett, 1977): the set of $m$ with $\sigma(m)<\infty$ has natural density $1$.
\item In the model where the parities are independent fair coins for ever, show that $\ln T^k(m)$ drifts by $\frac12\ln\frac34$ per step on average, and
predict how many steps a typical $m$ needs to reach $1$. (The notebook compares the prediction with the truth.)
\end{enumerate}
\end{exercise}

\begin{figure}[ht]
\centering
\fig{ch11-collatz}
\caption{The proportion of integers $2\le m<10^6$ with $\sigma(m)\le k$, that is with an iterate below $m$ within $k$ steps (blue), against the lower bound
$P(B_k<k\frac{\ln2}{\ln3})$ of Problem~\ref{pb:everett} (red). Both tend to $1$: almost every integer eventually goes down.}
\label{fig:collatz}
\end{figure}

\begin{exercise}[$\star\star\star$]\label{pb:holte}
\emph{The carries of a long addition.} Problem~\ref{pb:carries} studied the carries when two random numbers are added. Now add $n\ge2$
numbers written in base $b\ge2$, all digits independent and uniform on $\{0,\dots,b-1\}$. Let $C_0=0$ and let $C_i$ be the carry out of position
$i$, so that $C_i=\lfloor(C_{i-1}+D_{1,i}+\dots+D_{n,i})/b\rfloor$, where $D_{r,i}$ is the $i$-th digit of the $r$-th number.
\begin{enumerate}[label=(\alph*)]
\item Show that $(C_i)$ is a Markov chain on $\{0,\dots,n-1\}$. Call its transition matrix $P_b$.
\item Show that $P_b^m=P_{b^m}$ for every $m\ge1$.
\item Let $U_1,\dots,U_n$ be independent and uniform on $[0,1[$. Show that $|P_B(i,j)-P(\lfloor U_1+\dots+U_n\rfloor=j)|\le\frac{2n^2}B$ for all $i$, $j$ and $B$.
\item Deduce Holte's theorem: whatever the base, the stationary law of the carries is $\pi(j)=\frac{A(n,j)}{n!}$, the Eulerian numbers of
Problem~\ref{pb:irwinhall}, and $\|P_b^m(i,\cdot)-\pi\|\le\frac{2n^2}{b^m}$.
\item When three long decimal numbers are added, in what proportion of positions is the carry $1$? What is the average carry when $n$
numbers are added?
\end{enumerate}
\end{exercise}

\begin{exercise}[Journey]\label{pb:coprimewalk}
\emph{A drunkard among the visible points.} A point $(x,y)\in\Z^2$ is \emph{visible} from the origin if $\gcd(x,y)=1$; a uniformly chosen point of a
large square is visible with probability close to $\frac6{\pi^2}$. Let $(X_t,Y_t)$ be the simple random walk on $\Z^2$ started at $0$.
\begin{enumerate}[label=(\alph*)]
\item Show that $U_t=X_t+Y_t$ and $V_t=X_t-Y_t$ are two independent simple random walks on $\Z$.
\item Let $p$ be an odd prime. Show that $P(p\mid X_t\text{ and }p\mid Y_t)=P(p\mid U_t)^2$, and that $P(p\mid U_t)\to\frac1p$. Show that $X_t$ and $Y_t$ are never both even when $t$ is
odd, and are both even with probability tending to $\frac12$ when $t$ is even.
\item Show that $P(p\mid U_t)\le\frac1p+\frac2{\sqrt t}$ for every odd prime $p$ and every $t\ge1$. (Write the probability as an average of $\cos^t\frac{2\pi k}p$.)
\item Prove that $P\bigl(\gcd(X_t,Y_t)=1\bigr)$ tends to $\frac8{\pi^2}$ along odd $t$, and to $\frac4{\pi^2}$ along even $t$. (For finitely many primes, use the walk of $(U_t,V_t)$
modulo $4\prod p$ and Problem~\ref{pb:period}; for the others, use (c).)
\end{enumerate}
\end{exercise}

\section{Chains with a story}

The last problems are chains with a history: the range of a drunkard, a walk round a ring, three gamblers, the genetics of a
small population, two urns of Daniel Bernoulli, Kac's ring, and a card trick.

\begin{exercise}[$\star\star$]\label{pb:range}
\emph{How much ground does a drunkard cover?} Let $S$ be the simple random walk on $\Z$ started at $0$, and $R_t$ the number of distinct
points among $S_0,\dots,S_t$. Write $u_{2k}=\binom{2k}k4^{-k}$.
\begin{enumerate}[label=(\alph*)]
\item Show that $\mathbb E[R_t]=\sum_{s=0}^tP(S_1\ne0,\dots,S_s\ne0)$. (A new point is found at time $s$ when $S_s$ differs from all earlier positions; read the path backwards
from time $s$.)
\item Using Tome~3, Theorem~7.3, show that $\mathbb E[R_{2m}]=(4m+1)\,u_{2m}$.
\item Deduce that $\mathbb E[R_t]\sim\sqrt{8t/\pi}$.
\item In $\Z^d$ with $d\ge3$, show that $\frac1t\mathbb E[R_t]$ tends to the probability that the walk never returns to $0$.
\end{enumerate}
\end{exercise}

\begin{exercise}[$\star\star$]\label{pb:cyclecover}
\emph{The last vertex of the ring.} A walker moves on the cycle $\Z/n\Z$ ($n\ge3$), one step clockwise or anticlockwise with probability
$\frac12$ each, starting at $0$.
\begin{enumerate}[label=(\alph*)]
\item Show that the set of visited vertices is always an arc containing $0$, and that once it has $k<n$ vertices, the expected
time to add a new one is $k$.
\item Deduce that the expected time to visit every vertex is $\frac{n(n-1)}2$.
\item Show that the last vertex to be visited is uniformly distributed on $\{1,\dots,n-1\}$.
\item Show that the same is true for the walk on the complete graph $K_n$, but not for the walk on the star with three branches
started at the end of a branch.
\end{enumerate}
\end{exercise}

\begin{exercise}[$\star\star$]\label{pb:threegamblers}
\emph{Three gamblers.} Three players start with $a$, $b$, $c$ coins ($a,b,c\ge1$), and $s=a+b+c$. In each round two of the players who still
have coins are chosen uniformly, and a fair coin decides which of the two gives one coin to the other. A player with no coins
leaves the game. Let $X_t$, $Y_t$, $Z_t$ be the fortunes after $t$ rounds.
\begin{enumerate}[label=(\alph*)]
\item Show that $X_tY_tZ_t+\frac s3t$ is a martingale as long as all three players are in the game, and deduce that the expected number of
rounds before the first player leaves is $\frac{3abc}s$.
\item Show that $X_tY_t+Y_tZ_t+Z_tX_t+t$ is a martingale until one player holds all the coins, and that the expected duration of the whole
game is $ab+bc+ca$.
\item Show that the first player ends up with all the coins with probability $\frac as$.
\item With $4$ coins each, how long is the game on average, and how long before the first player leaves?
\end{enumerate}
\end{exercise}

\begin{exercise}[$\star\star\star$]\label{pb:moran}
\emph{A small population forgets.} A population has $N$ individuals of two types $A$ and $B$. At each step one individual, chosen
uniformly, reproduces, and one individual, chosen uniformly and independently (possibly the same one), dies and is replaced by the
offspring. Let $X_t$ be the number of individuals of type $A$, with $X_0=i$. (This is the Moran model of population genetics.)
\begin{enumerate}[label=(\alph*)]
\item Show that $X$ is a martingale, and that type $A$ takes over the whole population with probability $\frac iN$.
\item Show that $\mathbb E[X_{t+1}(N-X_{t+1})\mid X_t]=(1-\frac2{N^2})X_t(N-X_t)$: the diversity decays geometrically.
\item Show that the expected time until one type has disappeared is
\[
t_i=Ni\sum_{x=i}^{N-1}\frac1x+N(N-i)\sum_{x=N-i+1}^{N-1}\frac1x.
\]
(The walk moves up or down with equal probability when it moves; count the expected number of moves made from each state.)
\item Deduce that a single mutant ($i=1$) is resolved in time $NH_{N-1}\sim N\ln N$, while for $i=\lfloor pN\rfloor$ with $0<p<1$,
\[
t_i\sim N^2\Bigl(p\ln\frac1p+(1-p)\ln\frac1{1-p}\Bigr).
\]
\end{enumerate}
\end{exercise}

\begin{exercise}[$\star\star\star$]\label{pb:bernoullilaplace}
\emph{Daniel Bernoulli's two urns.} Two urns contain $n$ balls each; altogether $n$ balls are red and $n$ are white. At each step a ball is
drawn uniformly from each urn, and the two balls are exchanged. Let $X_t$ be the number of red balls in the first urn.
\begin{enumerate}[label=(\alph*)]
\item Find the transition probabilities of $X$, and show that the chain is reversible with respect to the hypergeometric law
$\pi(x)=\binom nx^2\binom{2n}n^{-1}$.
\item Let $(Pf)(x)=\mathbb E[f(X_1)\mid X_0=x]$. Show that if $f$ is a polynomial of degree $k$, so is $Pf$, and that the coefficient of $x^k$ is multiplied by
\[
\lambda_k=1-\frac{k(2n-k+1)}{n^2}.
\]
Deduce that the eigenvalues of the transition matrix are $\lambda_0,\dots,\lambda_n$.
\item Start with all red balls in the first urn. Show that $\mathbb E[X_t]-\frac n2=\frac n2(1-\frac2n)^t$, and deduce that for $n\ge4$ and $t\le\frac n4\ln n$ this is at least
$\frac{\sqrt n}3$, while under $\pi$ the standard deviation of $X$ is less than $\frac{\sqrt n}2$.
\end{enumerate}
\end{exercise}

Figure~\ref{fig:kacring} shows one ring: Boltzmann and Zermelo are both right, at different time scales.

\begin{exercise}[$\star\star$]\label{pb:kacring}
\emph{Kac's ring.} Label the sites of a circle by $\Z/n\Z$, and the edge from $j$ to $j+1$ by $j$. Each edge is marked independently with
probability $\mu\in\left]0,1\right[$; put $\varepsilon_j=-1$ if edge $j$ is marked and $\varepsilon_j=1$ otherwise, and let $M$ be the number of marked edges. Each site carries a
ball of colour $c_i=\pm1$ ($+1$ for black). At each tick every ball moves from its site $j$ to $j+1$, and changes colour if edge $j$ is
marked. Let $\Delta_t$ be the sum of the colours (the number of black balls minus the number of white ones) after $t$ ticks.
\begin{enumerate}[label=(\alph*)]
\item Show that the dynamics can be run backwards, that $\Delta_n=(-1)^M\Delta_0$, and that after $2n$ ticks every ball is back on its site with
its initial colour.
\item Show that the ball that starts at $i$ has colour $c_i\varepsilon_i\varepsilon_{i+1}\cdots\varepsilon_{i+t-1}$ after $t$ ticks, and that $\mathbb E[\Delta_t]=(1-2\mu)^t\Delta_0$ for $0\le t\le n$.
\item Show that $\operatorname{Var}(\Delta_t)\le n(2t-1)$ for $1\le t\le\frac{n+1}2$. Deduce that, for fixed $t$ and initial proportion $\beta$ of black balls, the
proportion of black balls at time $t$ tends in probability to $\frac12+(1-2\mu)^t(\beta-\frac12)$ as $n\to\infty$.
\item Starting with all balls black, the proportion of black balls is close to $\frac12$ after a few dozen ticks, yet it is $0$ or $1$ at
tick $n$, and $1$ at tick $2n$. Explain why there is no contradiction, and answer the questions of the think box.
\end{enumerate}
\end{exercise}

\begin{figure}[ht]
\centering
\fig{ch11-kacring}
\caption{One Kac ring with $n=2000$ sites, each edge marked with probability $\mu=\frac1{10}$ (here $M=217$ marks), all balls black at the start; left, the
first $80$ ticks, right, the whole period. The quantity $\frac{\Delta_t}n$ follows its mean $(1-2\mu)^t$ (dashed) down to $0$ and then fluctuates by about $\pm0.1$. At $t=n$ every ball has changed colour $217$ times, so all
are white; at $t=2n$ all are black again (Problem~\ref{pb:kacring}).}
\label{fig:kacring}
\end{figure}

\begin{exercise}[$\star\star$]\label{pb:gilbreath}
\emph{Gilbreath's card trick.} A deck has $n$ cards, numbered $1,\dots,n$ from the top. A \emph{Gilbreath shuffle} deals some number $k$ of cards from
the top, one at a time, onto a pile (which reverses their order), and then riffles that pile and the rest of the deck together, in
any way.
\begin{enumerate}[label=(\alph*)]
\item Show that the possible results are exactly the arrangements in which, for every $j\le n$, the top $j$ cards have distinct residues
modulo $j$. Deduce that there are $2^{n-1}$ of them.
\item The deck is arranged so that the suits repeat in the order spades, hearts, clubs, diamonds. Show that after a Gilbreath
shuffle, each of the blocks of cards $1$--$4$, $5$--$8$, $9$--$12$, \dots\ contains one card of each suit. (With colours alternating, the pairs
$1$--$2$, $3$--$4$, \dots\ each contain a red and a black card: this is Gilbreath's principle.)
\item Show that if $k$ is binomial with parameters $n$ and $\frac12$ and the riffle is uniform among the $\binom nk$ ways of interleaving, the result is
uniform among the $2^{n-1}$ possible arrangements.
\end{enumerate}
\end{exercise}

\section{Hints}

\begin{hint}{pb:returnbound}
(a) $X_{2t}=X_t\cdot(X_t^{-1}X_{2t})$, and the second factor has the law of $X_t$, inverted by symmetry. Then Cauchy--Schwarz, and $\sum_gp_t(g)=1$. (b) $p_{t+1}$ is
an average of translates of $p_t$, and the $\ell^2$ norm is convex. (c) Split $2t+1=t+(t+1)$.
\end{hint}

\begin{hint}{pb:period}
(a) In a finite group the semigroup generated by $S$ is a group. (b) Work in $G/N$, where the walk moves deterministically. (c) The cyclic classes of
Tome~3, Lemma~8.4 are $N,sN,s^2N,\dots$; to see that $N$ is normal, write $s^{-1}$ as a product of elements of $S$. (d) Signs. (e) $S=\{2,3\}\subseteq2\langle3\cdot2^{-1}\rangle$.
\end{hint}

\begin{hint}{pb:rollingdie}
(a) Tome~2, Proposition~1.5 and Exercise~1.6: a quarter turn about a face axis is a $4$-cycle of the diagonals. (b) Colour the board like a
chessboard. (c) A rotation of order $3$ fixes a diagonal; two such rotations about different diagonals generate $A_4$. (d) Problem~\ref{pb:period}, and
the action of $A_4$ on the six faces. (e) The uniform law is stationary.
\end{hint}

\begin{hint}{pb:dodecaturns}
(a) $\sum_hf(h)\rho(h)$ commutes with every $\rho(g)$; Schur's lemma, then the orthogonality relations to expand the class function $\mu^{*t}$ in
characters. (b) The ratios $\frac{\chi(g)}{\chi(1)}$ on a $3$-cycle are $1,0,0,\frac14,-\frac15$. (c) On a $5$-cycle the characters of dimensions $4$ and $5$ take the values
$-1$ and $0$.
\end{hint}

\begin{hint}{pb:gapdiam}
(a) $P$ is a symmetric matrix. (b) For $x=s_1\cdots s_\ell$, $(f(gx)-f(g))^2\le\ell\sum_i(f(gs_1\cdots s_i)-f(gs_1\cdots s_{i-1}))^2$; then sum over $g$, and substitute $g'=gs_1\cdots s_{i-1}$.
(c) The diameters are $\lfloor\frac n2\rfloor$ and $\binom n2$.
\end{hint}

\begin{hint}{pb:treewalk}
(b) Tome~3, Corollary~7.10. (c) Compute $\frac14h(k-1)+\frac34h(k+1)$. (d) Catalan numbers, and $\binom{2m}m\ge\frac{4^m}{2m+1}$. (e) The same computations with $\frac1q$ and $\frac{q-1}q$.
\end{hint}

\begin{hint}{pb:commute}
(a) Count visits along the cycles of length $\tau$ and use the strong Markov property: $\nu(x)=\mathbb E_a[\text{visits to }x\text{ before }\tau]$ satisfies $\nu P=\nu$. (b) Each
excursion from $a$ reaches $b$ before coming back with probability $\frac1{\deg(a)R}$. (d) $R(a,b)\le1$ along an edge; a closed walk round a
spanning tree uses each tree edge once in each direction.
\end{hint}

\begin{hint}{pb:foster}
(a) $L$ is symmetric with kernel $\R\mathbf1$, so its image is $W$. Take $v=-\frac hI$. (b) $R(x,y)=\langle e_x-e_y,L^{-1}(e_x-e_y)\rangle$ on $W$, and $L=\sum_{\{x,y\}}(e_x-e_y)(e_x-e_y)^{\mathsf T}$. (c) An
automorphism preserves resistances; the exchange gives $\mathbb E_x[T_y]=\mathbb E_y[T_x]$.
\end{hint}

\begin{hint}{pb:antipode}
(a) If $e_k$ is the expected time to go from $k$ to $k-1$, then $P(k,k-1)e_k=1+P(k,k+1)e_{k+1}$; with detailed balance for the law $\pi(k)\propto n_k$,
induct from the far end. (b) Layer sizes from $b$: octahedron $1,4,1$; cube $1,3,3,1$;
icosahedron $1,5,5,1$; dodecahedron $1,3,6,6,3,1$. (c) Levels are Hamming weights. (d) By symmetry $\mathbb E_a[T_b]=\mathbb E_b[T_a]$.
\end{hint}

\begin{hint}{pb:rayleigh}
(a) For a unit flow $\theta$, write $\theta=i+\eta$ and show $\sum_{\{x,y\}}i(x,y)\eta(x,y)=0$ by summation by parts. (b) A flow on the smaller graph is a flow on
the larger one. (c) The flow through each $\Pi_k$ is at least $1$; Cauchy--Schwarz. (d) Cut between the boundaries of $B_k$ and $B_{k+1}$. (e) Tome~3,
Exercise~7.8(b) on $H\cap B_n$ with its boundary merged.
\end{hint}

\begin{hint}{pb:aldousbroder}
(b) Each step resets the arc of the vertex that is left. (c) For a given $(T',y)$, the predecessors are obtained by choosing a child $x$ of $y$ and a
neighbour $z$ of $y$ in the subtree of $x$. (d) The last-exit tree of the forward path is the first-entrance tree of the reversed path.
\end{hint}

\begin{hint}{pb:terras}
(a) Induction on $k$: $T$ maps the two classes $r$ and $r+2^k$ modulo $2^{k+1}$ to classes modulo $2^k$ that differ by $3^{\epsilon}2^{k-1}$. (c) Induction.
\end{hint}

\begin{hint}{pb:everett}
(a) $T^k(m)=\lambda m+\rho$ with $\lambda<1$. (b) Problem~\ref{pb:terras}(a): the vectors with few ones correspond to that many residue classes. (c) $\frac{\ln2}{\ln3}>\frac12$ and the
weak law of large numbers. (d) Each step multiplies by $\frac12$ or roughly $\frac32$.
\end{hint}

\begin{hint}{pb:holte}
(a) The new carry depends on the old one and on fresh digits. (b) Add $m$ digits at a time. (c) Couple: $D_r=\lfloor BU_r\rfloor$. (d) A stationary law of
$P_b$ is stationary for $P_{b^m}$; let $m\to\infty$. (e) Symmetry $U\mapsto1-U$.
\end{hint}

\begin{hint}{pb:coprimewalk}
(a) Each step changes $U$ and $V$ by $\pm1$, with the four sign pairs equally likely. (b) $X=\frac{U+V}2$ and $Y=\frac{U-V}2$. (c) $|\cos y|\le e^{-2y^2/\pi^2}$ for $|y|\le\frac\pi2$, and
a Gaussian integral. (d) Truncate at primes $\le P$ and let $P\to\infty$ after $t$.
\end{hint}

\begin{hint}{pb:range}
(a) Reverse the steps: $S_s\ne S_r$ for all $r<s$ iff the reversed walk does not return to $0$ in $s$ steps. (b) Returns only happen at even times, and
$\sum_{k=0}^mu_{2k}=(2m+1)u_{2m}$. (c) $u_{2m}\sim\frac1{\sqrt{\pi m}}$. (d) Cesàro.
\end{hint}

\begin{hint}{pb:cyclecover}
(a) Gambler's ruin on an interval of length $k+1$. (c) For $v$ to be last, the walk must reach one neighbour of $v$ and then go all the way round
to the other neighbour without touching $v$. (d) On the star, where can the walk be last?
\end{hint}

\begin{hint}{pb:threegamblers}
(a) Compute the conditional expectation of the change for each pair, then average. Optional stopping needs $\mathbb E[\tau]<\infty$: Tome~3, Lemma~7.8 on a finite
state space. (c) $X_t$ is a bounded martingale.
\end{hint}

\begin{hint}{pb:moran}
(b) $\mathbb E[(X_{t+1}-X_t)^2\mid X_t]=\frac{2X_t(N-X_t)}{N^2}$. (c) From $x$, a move happens with probability $q_x=\frac{2x(N-x)}{N^2}$, so each move from $x$ costs $\frac1{q_x}$ steps on
average; the simple walk on $\{0,\dots,N\}$ started at $i$ makes on average $\frac{2i(N-x)}N$ moves from $x\ge i$, and $\frac{2x(N-i)}N$ from $x\le i$. (d) Riemann sums.
\end{hint}

\begin{hint}{pb:bernoullilaplace}
(a) Detailed balance. (b) Expand $(x+\delta)^k$ where $\delta\in\{-1,0,1\}$ is the change: only $\mathbb E[\delta]$ and $\mathbb E[\delta^2]$ reach degree $k$. (c) $\ln(1-y)\ge-y-y^2$ for $0\le y\le\frac12$.
\end{hint}

\begin{hint}{pb:kacring}
(b) The edges $i,\dots,i+t-1$ are distinct when $t\le n$. (c) Two balls whose windows of edges are disjoint have independent colours. (d) For which
$t$ does (c) say something?
\end{hint}

\begin{hint}{pb:gilbreath}
(a) Show that both conditions say that the top $j$ cards form a block of consecutive integers, for every $j$. For the residues, look first at the
bottom card, using $j=n-1$. For the count, build the arrangement from the bottom. (b) Among the top $4m$ cards each residue modulo $4$ appears $m$ times. (c) Count the pairs ($k$, interleaving)
that give each arrangement.
\end{hint}

\section{Solutions}

\begin{solution}{pb:returnbound}
\idea Two halves of a walk of length $2t$ are independent copies of a walk of length $t$, and symmetry turns their meeting into a scalar
product. Everything follows from Cauchy--Schwarz.

\writeup (a) Write $X_{2t}=X_tW$ with $W=\xi_{t+1}\cdots\xi_{2t}$, which is independent of $X_t$ and has the law $p_t$. Since $\mu$ is symmetric,
$X_t^{-1}=\xi_t^{-1}\cdots\xi_1^{-1}$ has the law of $X_t$, so $p_t(h^{-1})=p_t(h)$. Hence
\[
p_{2t}(g)=\sum_hp_t(h)\,p_t(h^{-1}g)=\sum_hp_t(h)\,p_t(g^{-1}h).
\]
For $g=e$ this is $\sum_hp_t(h)^2$. For general $g$, Cauchy--Schwarz gives
\[
p_{2t}(g)\le\Bigl(\sum_hp_t(h)^2\Bigr)^{1/2}\Bigl(\sum_hp_t(g^{-1}h)^2\Bigr)^{1/2}=p_{2t}(e),
\]
since $h\mapsto g^{-1}h$ is a bijection. Finally $1=(\sum_hp_t(h))^2\le|G|\sum_hp_t(h)^2$.

(b) Write $\|f\|=(\sum_hf(h)^2)^{1/2}$. Since $p_{t+1}(g)=\sum_s\mu(s)p_t(gs^{-1})$ and each translate $g\mapsto p_t(gs^{-1})$ has norm $\|p_t\|$, the triangle inequality gives
$\|p_{t+1}\|\le\sum_s\mu(s)\|p_t\|=\|p_t\|$. By (a), $p_{2t+2}(e)=\|p_{t+1}\|^2\le\|p_t\|^2=p_{2t}(e)$.

(c) On $\Z/2\Z$ with $\mu$ the point mass at $1$ (which is symmetric), $p_{2t+1}(0)=0$. In general, as in (a),
$p_{2t+1}(g)=\sum_hp_t(h)p_{t+1}(h^{-1}g)\le\|p_t\|\,\|p_{t+1}\|\le\|p_t\|^2=p_{2t}(e)$, by (b).

(d) The formula comes from Fourier analysis on $\Z/n\Z$ (Tome~2, Theorem~11.14): the step has transform $\cos\frac{2\pi k}n$. The term $k=0$ is $\frac1n$ and the
others are non-negative, so $p_{2t}(0)\ge\frac1n$; each term decreases with $t$ since $\cos^2\le1$. More generally
$p_{2t}(j)=\frac1n\sum_k\cos^{2t}\frac{2\pi k}n\cos\frac{2\pi kj}n\le p_{2t}(0)$ term by term.

\tools{Tome~2, Chapter~1 (groups) and Theorem~11.14 (Fourier analysis on $\Z/n\Z$); Tome~3, Chapter~8; Cauchy--Schwarz (Tome~0).}
\leads The same computation gives $\sum_g\bigl(p_t(g)-\frac1{|G|}\bigr)^2=p_{2t}(e)-\frac1{|G|}$: the $\ell^2$ distance to the uniform law is read off the
probability of being back at the start. This identity is the entry point of the $\ell^2$ methods in Diaconis's book \emph{Group
Representations in Probability and Statistics} (1988), and it powers Problem~\ref{pb:dodecaturns}.
\end{solution}

\begin{solution}{pb:period}
\idea Modulo a normal subgroup containing all the ratios of steps, the walk moves like a clock hand. Conversely the cyclic classes of
a periodic walk are the cosets of a normal subgroup.

\writeup (a) Let $\Sigma$ be the set of products of one or more elements of $S$. It is closed under products. In a finite group each $s$ has finite
order $o(s)$, so $e=s^{o(s)}$ and $s^{-1}=s^{o(s)-1}$ lie in $\Sigma\cup\{e\}$. So $\Sigma\cup\{e\}$ is a subgroup containing $S$, hence equal to $G$. Every $g$ is reachable from $e$,
and from $x$ the walk reaches $xg$ along the same steps: the walk is irreducible.

(b) Every step lies in $gN$, and $N$ is normal, so $X_t=\xi_1\cdots\xi_t\in g^tN$. The images of $S$ in $G/N$ all equal $gN$ and generate $G/N$, so $G/N$ is
cyclic, generated by $gN$, of order $m=[G:N]\ge2$. The walk can be at $e$ at time $t$ only if $g^t\in N$, that is $m\mid t$: the period is a multiple of $m$. In
particular $P(X_t\in N)=0$ when $m\nmid t$, while the uniform law gives $N$ the mass $\frac1m$; the walk does not converge.

(c) Let $N_r$ be the set of $x$ that the walk reaches at some time $t\equiv r\pmod d$. Since $e$ is reachable from every $x$, say in $u$ steps, two times $t,t'$
at which $x$ is reached satisfy $t+u,t'+u\in d\Z$; so the sets $N_0,\dots,N_{d-1}$ are disjoint, and they cover $G$ by (a). If $x\in N_0$ is reached at time $t$ and $y\in N_0$
at time $t'$, then $xy$ is reached at time $t+t'$ (go to $x$, then copy the path to $y$); so $N=N_0$ is closed under products, and being finite and
containing $e$ it is a subgroup.

Fix $s\in S$. The walk goes from $s$ back to $e$ in some $u$ steps, so $s^{-1}$ is reached at time $u$ with $1+u\equiv0$. If $x\in N$ is reached at time $t$, then $s^{-1}xs$
is reached at time $u+t+1\equiv0$: so $s^{-1}Ns\subseteq N$, hence $g^{-1}Ng\subseteq N$ for every $g$, since $S$ generates $G$ as a semigroup; $N$ is normal. If $x\in N_1$ is
reached at time $t\equiv1$, then $s^{-1}x$ is reached at time $u+t\equiv0$, so $x\in sN$; conversely $sN\subseteq N_1$. Thus $S\subseteq N_1=sN$, and likewise $N_r=s^rN$, so the $d$
classes are the cosets of $N$ and $[G:N]=d$.

If $S$ lies in no coset of a proper normal subgroup, then $N=G$ and $d=1$: the walk is irreducible and aperiodic, the uniform law is
stationary, and Tome~3, Theorem~8.11 gives convergence. If $S$ does lie in such a coset, (b) forbids convergence.

(d) All transpositions are odd, so $S$ lies in the coset $\tau A_n$ of the normal subgroup $A_n$ of index $2$. By (b) the period is even; since $\tau\tau=e$ the walk can
return at time $2$, so the period is $2$. The deck never converges: its sign alternates.

(e) The group $(\Z/p\Z)^*$ is abelian, so every subgroup is normal, and $\{2,3\}\subseteq gN$ means $2N=3N$, that is $q=3\cdot2^{-1}\in N$. The smallest such subgroup is $H=\langle q\rangle$. So
$S$ lies in a coset of a proper subgroup if and only if $q$ is not a primitive root, and then, by (b) and (c), the subgroup $N$ of (c) contains $H$
while the period is a multiple of $[G:H]$, so the period is exactly $[G:H]$. For $p=5$: $2^{-1}=3$ and $q=9\equiv4$, of order $2$ in a group of order $4$: period $2$
(indeed $2$ and $3$ are both non-squares modulo $5$, and the walk alternates between squares and non-squares). For $p=13$: $2^{-1}=7$, $q=21\equiv8$, and
$8^2\equiv12$, $8^4\equiv1$: $q$ has order $4$ in a group of order $12$, so the period is $3$. In both cases $2$ is a primitive root, so the hypothesis holds.

\tools{Tome~2, Chapters~1--2 (cosets, normal subgroups, quotients); Tome~3, Lemma~8.4 and Theorem~8.11.}
\leads The criterion is classical; Diaconis's book \emph{Group Representations in Probability and Statistics} (1988) states it in this form. How often
$3\cdot2^{-1}$ is a primitive root is a case of Artin's conjecture on primitive roots: Artin predicted in 1927 that $2$, for instance, is a
primitive root for about $37.4\%$ of all primes, and Hooley proved it in 1967 assuming the generalised Riemann hypothesis.
\end{solution}

\begin{solution}{pb:rollingdie}
\idea Tipping is a quarter turn, which is an odd permutation of the four diagonals, and a closed tour on a chessboard has even length. So
closed tours give even permutations only, and two small loops already give all of them.

\writeup Put the $x$-axis east, the $y$-axis north and the $z$-axis up, and describe the orientation of the cube by the rotation $R$ carrying
it from its initial position. A move replaces $R$ by $QR$, where $Q$ is one of the quarter turns
\[
R_E=\begin{pmatrix}0&0&1\\0&1&0\\-1&0&0\end{pmatrix},\quad R_N=\begin{pmatrix}1&0&0\\0&0&1\\0&-1&0\end{pmatrix},\quad R_W=R_E^{\mathsf T},\quad R_S=R_N^{\mathsf T}
\]
(rolling east sends the top face to the east, so $R_Ee_z=e_x$; similarly $R_Ne_z=e_y$).

(a) The rotation group of the cube permutes its four long diagonals, and this identifies it with $\mathfrak S_4$ (Tome~2, Proposition~1.5 and Exercise~1.6). A
quarter turn about the $z$-axis, $(x,y,z)\mapsto(-y,x,z)$, sends the diagonals through $(1,1,1)$, $(-1,1,1)$, $(-1,-1,1)$, $(1,-1,1)$ each to the next: it is a
$4$-cycle, an odd permutation (Tome~2, Corollary~4.3). Quarter turns about the other face axes are conjugate to it, hence also odd.

(b) Colour the board like a chessboard. Each move changes the colour of the square, so a tour that returns to its square has an even
number of moves, and its total rotation is a product of an even number of odd permutations: it is even. A quarter turn about the vertical
axis is a $4$-cycle, odd: impossible. A half turn is the square of a $4$-cycle, a product of two transpositions, even; by (c) it is possible.

(c) Multiplying the four matrices,
\[
R_SR_WR_NR_E=\begin{pmatrix}0&1&0\\0&0&-1\\-1&0&0\end{pmatrix},\qquad R_NR_WR_SR_E=\begin{pmatrix}0&-1&0\\0&0&1\\-1&0&0\end{pmatrix}.
\]
Both have trace $0=1+2\cos\theta$, so they are rotations by $\theta=\frac{2\pi}3$; their axes, the fixed vectors, are $(1,1,-1)$ and $(1,-1,-1)$: two different long diagonals.
As permutations of the diagonals they are $3$-cycles fixing different diagonals, so they generate a subgroup of $A_4$ with two different
subgroups of order $3$, of order divisible by $3$ and larger than $3$. As $A_4$ has no subgroup of order $6$ (Tome~2, Proposition~4.6), the subgroup is $A_4$
itself. Concatenating closed tours multiplies their rotations, so the rotations of closed tours form a group, which contains $A_4$ and is
contained in $A_4$ by (b).

(d) The orientation performs a random walk on the rotation group (with steps multiplied on the left, which changes nothing in
Problem~\ref{pb:period}). Its steps lie in the odd coset of the normal subgroup $A_4$, so the walk has period $2$ and does not converge. At even
times it is a walk on $A_4$ whose steps (two moves) include the identity ($R_WR_E=I$) and generate $A_4$ by (c); it converges to the uniform law on
$A_4$, and at odd times $X_{2t+1}$ converges to the uniform law on the odd coset. The body face on top is $R^{-1}e_z$. The group $A_4$ acts transitively on
the six faces $\pm e_x,\pm e_y,\pm e_z$: the rotations of order $3$ about $(1,1,1)$ permute $e_x,e_y,e_z$ cyclically, and the half turn about the $z$-axis sends $e_x$
to $-e_x$. By orbit--stabiliser (Tome~2, Theorem~3.5) every face is hit by exactly $2$ of the $12$ elements, so under the uniform law on $A_4$ (and also on
$A_4g$, which moves $g^{-1}e_z$ by a uniform element of $A_4$) each face is on top with probability $\frac16$. Both parities give $\frac16$, so the face on top converges.

(e) The orientation walk is irreducible on $24$ states (its steps and their products generate the whole group) and its stationary law is
uniform, so the expected return time is $24$ (Tome~3, Theorem~8.8).

\tools{Tome~2, Proposition~1.5, Exercise~1.6, Corollary~4.3, Proposition~4.6 and Theorem~3.5; Tome~3, Theorem~8.8; Problem~\ref{pb:period}.}
\leads The same parity count is the reason why half of the positions of the fifteen puzzle cannot be reached (Tome~2, Theorem~4.7).
The pair (square, orientation) is itself a random walk on a group, in which the translations of $\Z^2$ are twisted by the rotations of
the cube; the notebook estimates how often the rolling die is back on its square in its initial orientation.
\end{solution}

\begin{solution}{pb:dodecaturns}
\idea A law that is constant on conjugacy classes is diagonalised by the characters: each irreducible representation turns convolution into
multiplication by a number. For $A_5$ and $3$-cycles the numbers are $1,0,0,\frac14,-\frac15$.

\writeup (a) Let $\rho$ be irreducible of dimension $d$ and $A=\sum_hf(h)\rho(h)$ with $f$ a class function. Then $\rho(g)A\rho(g)^{-1}=\sum_hf(h)\rho(ghg^{-1})=A$ after the
substitution $h\mapsto g^{-1}hg$, so $A$ is a scalar $cI$ by Schur's lemma (Tome~2, Lemma~11.4); taking traces, $cd=\hat f(\chi)$. For two functions, with
$(f*g)(x)=\sum_hf(h)g(h^{-1}x)$, one checks $\sum_x(f*g)(x)\rho(x)=\bigl(\sum_hf(h)\rho(h)\bigr)\bigl(\sum_yg(y)\rho(y)\bigr)$. The law of $X_t=\xi_1\cdots\xi_t$ is $\mu^{*t}$, which is again a
class function, so $\widehat{\mu^{*t}}(\chi)=d\,(\hat\mu(\chi)/d)^t$.

The irreducible characters form an orthonormal basis of the class functions (Tome~2, Theorems~11.8 and~11.10). For a class function $F$,
$\langle F,\psi\rangle=\frac1{|G|}\sum_hF(h)\overline{\psi(h)}=\frac1{|G|}\hat F(\bar\psi)$, and $\bar\psi$ is again an irreducible character. Hence
$F=\frac1{|G|}\sum_\chi\hat F(\chi)\bar\chi$, and with $F=\mu^{*t}$ this is the formula.

(b) The character table of $A_5$, on the classes of $e$, $(12)(34)$, $(123)$, $(12345)$, $(13524)$, has rows $(1,1,1,1,1)$, $(3,-1,0,\varphi,1-\varphi)$,
$(3,-1,0,1-\varphi,\varphi)$, $(4,0,1,-1,-1)$, $(5,1,-1,0,0)$, with $\varphi=\frac{1+\sqrt5}2$. For $\mu$ uniform on the $3$-cycles, $\frac{\hat\mu(\chi)}{\chi(1)}=\frac{\chi((123))}{\chi(1)}$ takes the values $1,0,0,\frac14,-\frac15$. For $g=e$
and $t\ge1$ the formula gives
\[
\mu^{*t}(e)=\frac1{60}\Bigl(1+16\cdot4^{-t}+25\Bigl(-\frac15\Bigr)^t\Bigr).
\]
For $t=1$ this is $0$, as it should be, and for $t=2$ it is $\frac3{60}$: of the $400$ pairs of turns, $20$ undo each other.

(c) The rotations by $\pm\frac{2\pi}5$ about the six vertex axes form one conjugacy class of $12$ elements of order $5$. At such a $g$ the characters of
dimensions $4$ and $5$ are $-1$ and $0$, and the $3$-dimensional ones have ratio $0$, so $\mu^{*t}(g)=\frac1{60}(1-4\cdot4^{-t})$ and the probability asked for is
$\frac15(1-4^{1-t})$, for instance $\frac3{20}$ after two turns. In general $|\mu^{*t}(g)-\frac1{60}|\le\frac1{60}(16\cdot4^{-t}+25\cdot5^{-t})\le\frac{41}{60}4^{-t}$ for every $g$, so the distance to the uniform
law is at most $\frac{41}2\,4^{-t}$; it is not smaller order, since at $e$ the deviation is about $\frac{16}{60}4^{-t}$.

\tools{Tome~2, Lemma~11.4, Theorems~11.8 and 11.10, Exercise~11.7 (the table of $A_5$), Theorem~12.5.}
\leads Diaconis and Shahshahani used exactly this method in 1981, with the characters of $\mathfrak S_n$, to show that random transpositions mix a deck
after about $\frac12n\ln n$ steps and not before, one of the first cutoff results. Tome~3, Chapter~8 proves the upper bound by Broder's
probabilistic argument instead.
\end{solution}

\begin{solution}{pb:gapdiam}
\idea A function that varies a lot overall must vary along some edge: connect any two elements by a short path and pay for its
length twice, once in Cauchy--Schwarz and once in the number of paths through an edge.

\writeup (a) Expanding the square and using that $g\mapsto gs$ is a bijection,
$\sum_{g,s}(f(gs)-f(g))^2=2|S|\sum_gf(g)^2-2\sum_gf(g)\sum_sf(gs)$, so $\mathcal E(f)=\frac1{|G|}\sum_gf(g)\bigl(f(g)-(Pf)(g)\bigr)=\langle f,(I-P)f\rangle$. Similarly
$\operatorname{Var}(f)=\frac1{|G|}\sum f^2-(\frac1{|G|}\sum f)^2$. Since $S=S^{-1}$, $P$ is a symmetric matrix, hence self-adjoint for $\langle\cdot,\cdot\rangle$, with an orthonormal basis of eigenvectors, the
constant $\mathbf 1$ for $\lambda_1=1$. Neither $\mathcal E$ nor $\operatorname{Var}$ changes when a constant is added to $f$, so we may assume $f\perp\mathbf1$, and then $\frac{\mathcal E(f)}{\operatorname{Var}(f)}=\frac{\langle f,(I-P)f\rangle}{\langle f,f\rangle}\ge1-\lambda_2$,
with equality for an eigenvector of $\lambda_2$.

(b) For $x\in G$ fix a word $x=s_1\cdots s_\ell$ with $\ell\le d$, and put $w_i=s_1\cdots s_i$. Then $f(gx)-f(g)=\sum_{i=1}^\ell\bigl(f(gw_i)-f(gw_{i-1})\bigr)$, so by Cauchy--Schwarz
\[
\sum_g\bigl(f(gx)-f(g)\bigr)^2\le\ell\sum_{i=1}^\ell\sum_g\bigl(f(gw_{i-1}s_i)-f(gw_{i-1})\bigr)^2\le\ell^2\sum_{g'}\sum_{s\in S}\bigl(f(g's)-f(g')\bigr)^2,
\]
after the substitution $g'=gw_{i-1}$. The right side is at most $2d^2|G||S|\mathcal E(f)$. Summing over the $|G|$ values of $x$ (with $h=gx$),
$\operatorname{Var}(f)=\frac1{2|G|^2}\sum_x\sum_g(f(gx)-f(g))^2\le|S|d^2\mathcal E(f)$. With (a), $1-\lambda_2\ge\frac1{|S|d^2}$.

(c) For $\Z/n\Z$, $|S|=2$ and $d=\lfloor\frac n2\rfloor$, so $1-\lambda_2\ge\frac2{n^2}$, while the truth is $1-\cos\frac{2\pi}n\sim\frac{2\pi^2}{n^2}$: only the constant is lost. For $\mathfrak S_n$
with adjacent transpositions, $|S|=n-1$ and $d=\binom n2$ (an adjacent transposition changes the number of inversions by one, and bubble
sort shows that $\binom n2$ steps suffice), so $1-\lambda_2\ge\frac4{(n-1)^3n^2}$. The true gap is $\frac2{n-1}(1-\cos\frac\pi n)\sim\frac{\pi^2}{n^3}$ (the notebook checks it for $n\le6$): here the crude
path argument loses a factor of order $n^2$.

\tools{Tome~3, Proposition~8.16 (eigenvalues of reversible chains); spectral theorem for symmetric matrices (Tome~0); Cauchy--Schwarz.}
\leads Choosing the paths more cleverly, so that few of them use each edge, is the method of canonical paths of Jerrum and Sinclair and
the comparison theorems of Diaconis and Saloff-Coste (1993), which transfer a gap from a walk one understands to one that one
does not. For adjacent transpositions, Wilson (2004) showed that the mixing time is of order $n^3\ln n$.
\end{solution}

\begin{solution}{pb:treewalk}
\idea On the tree the distance to the root is a walk with a drift away from it, so returns are rare: the return probability decays
exponentially, at the rate $\frac{\sqrt3}2$, while on $\Z^2$ it decays only like $\frac1{\pi n}$.

\writeup (a) A reduced word of length $k\ge1$ ends with some letter $\ell$; multiplying by $\ell^{-1}$ shortens it, and by any of the three other letters
lengthens it, without cancellation. From $e$ every letter gives length $1$. The transition probabilities depend only on the length, so $(D_t)$ is
a Markov chain with the stated transitions.

(b) The walk returns to $e$ exactly when $D$ returns to $0$, which it does if and only if, from $1$, it ever reaches $0$. For the walk on $\Z$ with
up-probability $p=\frac34$ and down-probability $q=\frac14$ started at $1$, this happens with probability $\frac qp=\frac13$ (Tome~3, Corollary~7.10). The number of visits
to $e$ is geometric: $P(V\ge k)=3^{-(k-1)}$, so $\mathbb E[V]=\frac1{1-1/3}=\frac32$.

(c) For $k\ge1$, $\frac14h(k-1)+\frac34h(k+1)=h(k)\bigl(\frac{\sqrt3}4+\frac3{4\sqrt3}\bigr)=\frac{\sqrt3}2h(k)$, and for $k=0$, $h(1)=\frac1{\sqrt3}\le\frac{\sqrt3}2$. So $\mathbb E[h(D_{t+1})]\le\frac{\sqrt3}2\mathbb E[h(D_t)]$, and by induction
$\mathbb E[h(D_{2n})]\le(\frac{\sqrt3}2)^{2n}$. Since $h\ge0$ and $h(0)=1$, $P(D_{2n}=0)\le\mathbb E[h(D_{2n})]$.

(d) Consider the paths with $D_1=1$, $D_s\ge1$ for $1\le s\le2n-1$, $D_{2n-1}=1$ and $D_{2n}=0$. Between times $1$ and $2n-1$ they make $n-1$ up-steps and $n-1$ down-steps
without going below $1$, in $C_{n-1}$ ways ($C_m=\frac1{m+1}\binom{2m}m$, the Catalan number), each of probability $(\frac34)^{n-1}(\frac14)^{n-1}$, and the last step has probability
$\frac14$. With $\binom{2m}m\ge\frac{4^m}{2m+1}$ we get $C_{n-1}\ge\frac{4^{n-1}}{(2n-1)n}\ge\frac{4^{n-1}}{2n^2}$, hence
\[
P(D_{2n}=0)\ge\frac14\cdot\frac{4^{n-1}}{2n^2}\Bigl(\frac3{16}\Bigr)^{n-1}=\frac1{8n^2}\Bigl(\frac34\Bigr)^{n-1}.
\]
Taking $2n$-th roots, both bounds tend to
$\sqrt{3/4}=\frac{\sqrt3}2$.

(e) On the $q$-regular tree the distance goes up with probability $\frac{q-1}q$ and down with probability $\frac1q$ (from $k\ge1$). With $h(k)=(q-1)^{-k/2}$,
$\frac1qh(k-1)+\frac{q-1}qh(k+1)=\frac{2\sqrt{q-1}}qh(k)$, and $h(1)\le\frac{2\sqrt{q-1}}q$ since $q\ge2$; the path count gives $C_{n-1}(\frac{q-1}{q^2})^{n-1}\frac1q$, whose $2n$-th root tends to
$\frac{2\sqrt{q-1}}q$. On $\Z^2$, $P(S_{2n}=0)=u_{2n}^2\sim\frac1{\pi n}$ (two independent coordinates after a rotation by $45^\circ$, as in Problem~\ref{pb:coprimewalk}), whose $2n$-th
root tends to $1$.

\tools{Tome~3, Corollary~7.10 (gambler's ruin against an infinite bank) and Theorem~7.3; Catalan numbers (Tome~0).}
\leads Kesten proved in 1959 that the return probabilities of a symmetric walk on a group decay exponentially if and only if the group is
not \emph{amenable}, the property that fails for the free group and that underlies the Banach--Tarski paradox (Tome~7). The number
$2\sqrt{q-1}$, the spectral radius of the tree's adjacency operator, is the benchmark of the Ramanujan graphs of Tome~11, Appendix~A.
\end{solution}

\begin{solution}{pb:commute}
\idea A cycle that starts and ends at $a$ spends, on average, time at each vertex in proportion to the stationary law. Applied to the round
trip $a\to b\to a$, it equates two ways of counting visits to $a$: by the stationary law, and by the escape probability of the electrical
network.

\writeup (a) Let $\nu(x)=\mathbb E_a\bigl[\sum_{t<\tau}\mathbf 1_{X_t=x}\bigr]$, which is finite since $\nu(x)\le\mathbb E_a[\tau]$. For each $y$,
\[
\nu(y)=\mathbb E_a\Bigl[\sum_{t=1}^{\tau}\mathbf 1_{X_t=y}\Bigr]=\sum_{t\ge1}\sum_xP_a(X_{t-1}=x,\ \tau>t-1)P(x,y)=\sum_x\nu(x)P(x,y).
\]
The first equality holds because $X_0=X_\tau=a$, so shifting the window by one changes nothing; the second uses the Markov property at time $t-1$,
the event $\{\tau>t-1\}$ being determined by $X_0,\dots,X_{t-1}$. So $\nu$ is a non-negative invariant measure of total mass $\mathbb E_a[\tau]$. By uniqueness
(Tome~3, Theorem~8.8), $\nu=\mathbb E_a[\tau]\,\pi$.

(b) Here $\mathbb E_a[\tau]=\mathbb E_a[T_b]+\mathbb E_b[T_a]$ by the strong Markov property at $T_b$. The visits to $a$ before $\tau$ all happen before $T_b$: the walk starts at $a$,
makes excursions from $a$, and stops counting once an excursion reaches $b$. Each excursion reaches $b$ before returning to $a$ with
probability $\frac1{\deg(a)R(a,b)}$, independently (Tome~3, Exercise~7.8(b)), so the number of visits to $a$ is geometric with mean $\deg(a)R(a,b)$. By (a)
it equals $\pi(a)\mathbb E_a[\tau]=\frac{\deg a}{2|\mathcal E|}\mathbb E_a[\tau]$, and the identity follows.

(c) On a path $0,1,\dots,m$, $R=m$ (resistors in series) and $|\mathcal E|=m$, so the identity predicts a commute time $2m^2$. Indeed the walk on the path
started at $0$ is $|S|$ for the simple walk $S$ on $\Z$, so $T_m$ is the exit time of $S$ from $\{-m+1,\dots,m-1\}$, of mean $m\cdot m=m^2$ by the gambler's ruin formula
$\mathbb E_x[\tau]=x(N-x)$ (Tome~3, Theorem~7.9, with $N=2m$, $x=m$); by symmetry the trip back also takes $m^2$. On a cycle of length $n$, two neighbours are joined by
resistances $1$ and $n-1$ in parallel, so $R=\frac{n-1}n$ and $2|\mathcal E|R=2(n-1)$. Directly: unrolling the cycle, the walk from a neighbour of $b$ hits $b$ when it
leaves $\{1,\dots,n-1\}$ from $1$, which takes $1\cdot(n-1)$ on average; each one-way trip takes $n-1$.

(d) If $a\sim b$, the edge $ab$ alone carries a unit flow of energy $1$, so $R(a,b)\le1$ by Thomson's principle (Problem~\ref{pb:rayleigh}(a)). Then $\mathbb E_a[T_b]<\mathbb E_a[T_b]+\mathbb E_b[T_a]\le2|\mathcal E|$. Take a spanning tree and a closed walk
$v_0,v_1,\dots,v_{2n-2}=v_0$ round it that uses each of its $n-1$ edges once in each direction. Following the walk from $v_0$, wait each time until the
random walk hits the next vertex $v_{i+1}$; when the last one is reached, every vertex has been visited. The expected time is at most
$\sum_i\mathbb E_{v_i}[T_{v_{i+1}}]=\sum_{\text{tree edges}}(\mathbb E_x[T_y]+\mathbb E_y[T_x])\le(n-1)\cdot2|\mathcal E|$, where the tour is started at the walk's starting vertex $v_0$.

\tools{Tome~3, Theorem~8.8, Corollary~8.15, Exercise~7.8 and Theorem~7.9.}
\leads The commute time identity is due to Chandra, Raghavan, Ruzzo, Smolensky and Tiwari (1989), and the cover time bound to Aleliunas,
Karp, Lipton, Lovász and Rackoff (1979), who used it to decide connectivity in logarithmic space with a random walk. In 2005 Reingold
removed the randomness: undirected connectivity can be decided deterministically in logarithmic space.
\end{solution}

\begin{solution}{pb:foster}
\idea The resistances are the entries of the inverse of the Laplacian on the hyperplane $\mathbf 1^\perp$, and summing them over the edges rebuilds
the Laplacian inside a trace, whose value is the dimension $n-1$.

\writeup (a) $L$ is symmetric and $\langle x,Lx\rangle=\sum_{\{u,w\}\in\mathcal E}(x_u-x_w)^2$, which vanishes only for constant $x$ ($G$ is connected). So $\ker L=\R\mathbf 1$, and the image of $L$ is
$(\ker L)^\perp=W$. As $e_a-e_b\in W$, there is $v$ with $Lv=e_a-e_b$, unique up to constants. Let $h$ be the voltage of Tome~3, Exercise~7.8: $h(a)=0$,
$h(b)=1$, harmonic elsewhere, and $I=\sum_{y\sim a}h(y)=-(Lh)(a)$. Then $(Lh)(x)=0$ for $x\ne a,b$, and $\sum_x(Lh)(x)=\langle\mathbf1,Lh\rangle=0$ gives $(Lh)(b)=I$. So $v=-\frac hI$
satisfies $Lv=e_a-e_b$, and $v(a)-v(b)=\frac1I=R(a,b)$.

(b) Let $M$ be the inverse of $L$ restricted to $W$. For $x\sim y$, $R(x,y)=\langle e_x-e_y,M(e_x-e_y)\rangle$, since $v=M(e_x-e_y)$ is a solution and $\langle e_x-e_y,v\rangle=v(x)-v(y)$. Write
$\langle u,Mu\rangle=\operatorname{tr}(M\,uu^{\mathsf T})$ for $u\in W$, and note that $L=\sum_{\{x,y\}\in\mathcal E}(e_x-e_y)(e_x-e_y)^{\mathsf T}$, as operators on $W$. Then
\[
\sum_{\{x,y\}\in\mathcal E}R(x,y)=\operatorname{tr}_W\Bigl(M\sum_{\{x,y\}}(e_x-e_y)(e_x-e_y)^{\mathsf T}\Bigr)=\operatorname{tr}_W(ML)=\dim W=n-1.
\]

(c) An automorphism $\gamma$ maps the voltage problem for $(a,b)$ to the one for $(\gamma a,\gamma b)$, so $R(\gamma a,\gamma b)=R(a,b)$. If the automorphisms act transitively
on edges, all edges have the same resistance, and Foster's theorem gives $|\mathcal E|R=n-1$. If some automorphism exchanges $x$ and $y$, it
exchanges the laws of the walks from $x$ and from $y$, so $\mathbb E_x[T_y]=\mathbb E_y[T_x]$, and the commute time identity gives $2\mathbb E_x[T_y]=2|\mathcal E|R=2(n-1)$.

(d) The cube, the Petersen graph and the dodecahedron have the required symmetry (their automorphism groups act transitively on
ordered pairs of neighbours). So the expected times are $n-1$: $7$, $9$ and $19$.

\tools{Tome~3, Exercise~7.8; Problem~\ref{pb:commute}; Tome~0 (symmetric matrices, trace); Tome~2, Chapter~3 (group actions).}
\leads Foster published the theorem in 1949. It has a one-line probabilistic meaning: for an edge $xy$, $R(x,y)$ is the probability that the
edge belongs to a uniformly chosen spanning tree (a consequence of Kirchhoff's work of 1847), and every spanning tree has $n-1$ edges.
Problem~\ref{pb:aldousbroder} shows how a random walk samples such a tree.
\end{solution}

\begin{solution}{pb:antipode}
\idea On a very symmetric polyhedron only the distance to the target matters, so the walk collapses to a walk on a few levels, and
a birth--death chain has explicit hitting times.

\writeup (a) Given the past, the walk moves to a uniform neighbour of $X_t$; if $X_t$ is at distance $k$ from $b$, the new distance is $k-1$, $k$ or $k+1$
with probabilities $\frac{c_k}r$, $\frac{a_k}r$, $\frac{b_k}r$, which depend only on $k$. So the distance is a Markov chain, a birth--death chain on $\{0,\dots,D\}$. Counting edges
between levels gives $n_kc_k=n_{k-1}b_{k-1}$, which is detailed balance for $\pi(k)\propto n_k$. Let $e_k$ be the expected time to go from $k$ to $k-1$. From $k$, one
step is made; then either the level is $k-1$, or it is $k$ again (cost $e_k$ more), or it is $k+1$ (cost $e_{k+1}+e_k$ more):
$e_k=1+\frac{a_k}re_k+\frac{b_k}r(e_{k+1}+e_k)$, that is $\frac{c_k}re_k=1+\frac{b_k}re_{k+1}$, with $b_D=0$. If $e_{k+1}=\frac r{n_{k+1}c_{k+1}}\sum_{i>k}n_i$, then $\frac{b_k}re_{k+1}=\frac{b_k}{n_{k+1}c_{k+1}}\sum_{i>k}n_i=\frac1{n_k}\sum_{i>k}n_i$
by detailed balance, so $e_k=\frac r{c_k}\bigl(1+\frac1{n_k}\sum_{i>k}n_i\bigr)=\frac r{n_kc_k}\sum_{i\ge k}n_i$. The formula holds for $k=D$ directly, and by induction downwards for all $k$.

(b) $\mathbb E_a[T_b]=\sum_{k=1}^De_k$. Octahedron ($r=4$, levels $1,4,1$, $c_1=1$, $c_2=4$): $e_2=\frac4{4}=1$, $e_1=\frac4{4}\cdot5=5$, total $6$. Cube ($r=3$, levels $1,3,3,1$, $c=1,2,3$):
$e_3=1$, $e_2=\frac3{6}\cdot4=2$, $e_1=\frac3{3}\cdot7=7$, total $10$. Icosahedron ($r=5$, levels $1,5,5,1$, $c=1,2,5$): $e_3=1$, $e_2=\frac5{10}\cdot6=3$, $e_1=\frac55\cdot11=11$, total $15$.
Dodecahedron ($r=3$, levels $1,3,6,6,3,1$, $c=1,1,1,2,3$): $e_5=1$, $e_4=\frac36\cdot4=2$, $e_3=\frac36\cdot10=5$, $e_2=\frac36\cdot16=8$, $e_1=\frac33\cdot19=19$, total $35$.

(c) Measure distances from $1\cdots1$ by the number of zeros, so the walk starts at level $m$. Level $k$ has $\binom mk$ vertices, each with $c_k=k$ neighbours at
level $k-1$. Hence $e_k=\frac m{k\binom mk}\sum_{i\ge k}\binom mi=\binom{m-1}{k-1}^{-1}\sum_{j\le m-k}\binom mj$. Summing over $k$ and substituting $k\mapsto m+1-k$ (with $\binom{m-1}{k-1}=\binom{m-1}{m-k}$) gives the
formula. In the stated form, the term $k=m$ is $2^m-1$ and the term $k=m-1$ is $\frac{2^m-1-m}{m-1}$; the terms $k=1$ and $k=2$ are $1$ and $\frac{m+1}{m-1}\le3$; and for
$3\le k\le m-2$ each term is at most $2^m\binom{m-1}2^{-1}$, so together they are at most $\frac{2m\,2^m}{(m-1)(m-2)}$. So the sum is $2^m(1+O(\frac1m))\sim2^m$.

(d) By symmetry (the antipodal map exchanges $a$ and $b$), $\mathbb E_a[T_b]=\mathbb E_b[T_a]$, so $R(a,b)=\frac{2\mathbb E_a[T_b]}{2|\mathcal E|}=\frac{\mathbb E_a[T_b]}{|\mathcal E|}$. Cube: $\frac{10}{12}=\frac56$ ohm; icosahedron:
$\frac{15}{30}=\frac12$ ohm; dodecahedron: $\frac{35}{30}=\frac76$ ohm.

\tools{Tome~3, Chapter~8 (detailed balance) and Problem~\ref{pb:commute}; Tome~2, Chapter~3 (orbits of the stabiliser).}
\leads Graphs in which the numbers $c_k,a_k,b_k$ are constant on levels around every vertex are the \emph{distance-regular} graphs, a central object of
algebraic combinatorics; their spectra are computed from the same tridiagonal matrix as the walk above (Brouwer, Cohen and Neumaier,
\emph{Distance-Regular Graphs}, 1989). The resistance of the cube, $\frac56$ ohm, is a classic puzzle of physics courses.
\end{solution}

\begin{solution}{pb:rayleigh}
\idea Current is the unit flow of least energy. Deleting edges shrinks the set of flows, so it can only increase the minimum; and any
family of disjoint cuts forces every flow to pay a toll at each cut. In the plane the tolls add up to a divergent harmonic series.

\writeup (a) $i(x,y)=v(x)-v(y)$ is antisymmetric, and $\sum_yi(x,y)=\deg(x)v(x)-\sum_{y\sim x}v(y)=(Lv)(x)$, which is $1$ at $a$, $-1$ at $b$ and $0$ elsewhere: a unit flow. Its
energy is $\sum_{\{x,y\}}(v(x)-v(y))^2=\langle v,Lv\rangle=v(a)-v(b)=R(a,b)$. If $\theta$ is another unit flow, $\eta=\theta-i$ has zero net flow at every vertex, and
\[
\sum_{\{x,y\}}i(x,y)\eta(x,y)=\frac12\sum_x\sum_{y\sim x}(v(x)-v(y))\eta(x,y)=\sum_xv(x)\sum_{y\sim x}\eta(x,y)=0,
\]
using antisymmetry of $\eta$ for the middle equality. Hence $\mathcal W(\theta)=\mathcal W(i)+\mathcal W(\eta)\ge R(a,b)$.

(b) A unit flow on the graph with fewer edges is a unit flow on the larger graph (put $0$ on the removed edges), with the same energy. The
minimum over the smaller set is larger.

(c) Let $\theta$ be a unit flow and $A$ the set of vertices that $a$ can reach without using an edge of $\Pi_k$, so $b\notin A$. Summing the net flows over $A$,
the contributions of edges inside $A$ cancel, and $1=\sum_{x\in A,\,y\notin A,\,x\sim y}\theta(x,y)$; all those edges belong to $\Pi_k$. So $\sum_{e\in\Pi_k}|\theta_e|\ge1$, and by Cauchy--Schwarz
$\sum_{e\in\Pi_k}\theta_e^2\ge\frac1{|\Pi_k|}$. The $\Pi_k$ are disjoint, so $\mathcal W(\theta)\ge\sum_k\frac1{|\Pi_k|}$; minimise over $\theta$.

(d) For $0\le k<n$ let $\Pi_k$ be the set of edges from $\partial B_k$ to $\partial B_{k+1}$ (with $\partial B_0=\{0\}$). There are $4(2k+1)$ of them (each side of the square $\partial B_k$ has
$2k+1$ points with one outward edge), they are disjoint, and every path from $0$ to $z_n$ crosses each of them. By (c),
\[
R(0,z_n)\ge\sum_{k=0}^{n-1}\frac1{4(2k+1)}\ge\frac14\int_0^n\frac{\dd x}{2x+1}=\frac18\ln(2n+1).
\]

(e) Let $H_n$ be the graph with the edges of $H$ inside $B_n$, the boundary merged into $z_n$. If $0$ is not connected to $z_n$ in $H_n$, the component of $0$ in $H$
stays inside $B_n$; it is finite, the walk on it is a finite irreducible chain, and it returns to $0$ for sure (Tome~3, Theorem~8.8). Otherwise, working in
the component of $0$ in $H_n$ and using Rayleigh, $R_{H_n}(0,z_n)\ge R(0,z_n)\ge\frac18\ln(2n+1)$, and by Tome~3, Exercise~7.8(b), the walk on $H$
started at $0$ reaches $\partial B_n$ before returning to $0$ with probability $\frac1{\deg_H(0)R_{H_n}(0,z_n)}\le\frac8{\ln(2n+1)}$ (until it reaches $\partial B_n$ the two walks coincide). Now
$P(\text{no return})\le P(\text{reach }\partial B_n\text{ before returning})+P(\text{stay in }B_n\text{ for ever without returning})$. The second term is $0$: from every vertex of $B_n$ that the walk
can reach there is a path in $H$ to $0$, which stays in $B_n$ or touches $\partial B_n$, so within $|B_n|$ steps the walk hits $\{0\}\cup\partial B_n$ with probability at least some
$\delta>0$, and it does so eventually with probability $1$ (as in Tome~3, Lemma~7.8). Letting $n\to\infty$, $P(\text{no return})=0$.

\tools{Tome~3, Exercise~7.8, Lemma~7.8, Theorem~8.8; Problem~\ref{pb:foster}(a); Cauchy--Schwarz.}
\leads The Nash-Williams criterion (1959) and Thomson's principle make recurrence a question about electrical networks, and Lyons and
Peres's book \emph{Probability on Trees and Networks} (2016) develops the theory. The same cutsets in $\Z^3$ have about $24k^2$ edges, so the
lower bound converges, and indeed one can build a unit flow of finite energy to infinity: the walk on $\Z^3$ is transient, as Pólya
proved (Tome~3, Theorem~7.13). The monotonicity is not obvious from the walk: removing edges changes all the transition probabilities.
\end{solution}

\begin{solution}{pb:aldousbroder}
\idea Follow the tree of last exits along the walk: it is a Markov chain on rooted trees whose stationary law weighs each tree by the
product of its transition probabilities. Reversing time turns last exits into first entrances, and the root's weight cancels.

\writeup (a) The first visit to each $v\ne r$ comes after the first visit to some neighbour $u$, which was visited earlier, so the kept edge joins $v$ to an
earlier vertex. Each $v\ne r$ has one kept edge to an earlier vertex: there are $n-1$ edges and following them always decreases the time of first
visit, so they connect every vertex to $r$ without cycles. They form a spanning tree.

(b) In $T$, rooted at $x$, the vertex $y$ has one outgoing arc, and its subtree (the vertices whose path to $x$ passes through $y$) does not contain
$x$. Deleting $y$'s arc separates this subtree, now rooted at $y$, from the rest, rooted at $x$; adding $x\to y$ attaches the rest below $y$. The result is
an arborescence rooted at $y$. Each step sets the arc of the vertex being left to the vertex entered next, and later steps change it only
when that vertex is left again. So after a path that visits every vertex, the arc of every vertex other than the endpoint points to the
vertex that followed its last visit, whatever the starting state.

(c) Fix $(T',y)$. A predecessor $(T,x)$ must have $x\to y$ in $T'$, so $x$ is a child of $y$ in $T'$, and $T=T'-(x\to y)+(y\to z)$ for some $z$, which must lie in the
subtree of $x$ in $T'$ (so that $T$ is an arborescence rooted at $x$); conversely each such choice of $(x,z)$ with $z\sim y$ gives a predecessor. Then
$w(T,x)=w(T',y)\frac{\deg x}{\deg y}$ (the arc of $y$ is new, the arc of $x$ is gone), and the transition has probability $\frac1{\deg x}$. So
\[
\sum_{(T,x)\to(T',y)}w(T,x)\frac1{\deg x}=w(T',y)\sum_{x}\sum_{z}\frac1{\deg y}=w(T',y),
\]
since every neighbour $z$ of $y$ lies in the subtree of exactly one child $x$ of $y$. So $w$, hence $\mu$, is stationary. Next $w(T,x)=\frac{\deg x}{\prod_u\deg u}$, so
$Z=\sum_x\sum_Tw(T,x)=\frac{2|\mathcal E|\tau(G)}{\prod_u\deg u}$ (each spanning tree has one orientation towards each root) and $\mu(T,x)=\frac{\deg x}{2|\mathcal E|\tau(G)}$. Summing over $T$, the root has law
$\frac{\deg x}{2|\mathcal E|}=\pi(x)$.

(d) Let $(T_0,X_0)$ have law $\mu$ and run the tree chain along the walk $X_0,\dots,X_N$, which is the simple random walk started from $\pi$. Then $(T_N,X_N)$ has
law $\mu$. Let $A_N$ be the event that the walk visits every vertex by time $N$; on $A_N$, $T_N$ is the last-exit tree of $X_0,\dots,X_N$, rooted at $X_N$, by (b).
The reversed path $Y_s=X_{N-s}$ is again a simple random walk started from $\pi$ (Tome~3, Proposition~8.14), and the last visit of $X$ to $v$ before $N$
is the first visit of $Y$ to $v$; the vertex following it in $X$ precedes it in $Y$. So, as undirected trees, the last-exit tree of $X$ is the
first-entrance tree $F$ of $Y$, and $X_N=Y_0$. Hence for every tree $T$ and vertex $r$,
\[
\bigl|P(F=T,\ Y_0=r,\ A_N)-\mu(T,r)\bigr|\le P(A_N^c).
\]
As $N\to\infty$, $P(A_N^c)\to0$ (the cover time is finite almost surely, Problem~\ref{pb:commute}(d)), and on $A_N$ the first-entrance tree of $Y_0,\dots,Y_N$ is that of the
whole walk $Y$. So $P(F=T,\,Y_0=r)=\mu(T,r)$ and $P(F=T\mid Y_0=r)=\frac{\mu(T,r)}{\pi(r)}=\frac1{\tau(G)}$.

\tools{Tome~3, Proposition~8.14 (reversibility) and Corollary~8.15; Problem~\ref{pb:commute}(d).}
\leads The algorithm was found independently by Broder and by Aldous around 1989--1990; part (c) is the Markov chain tree theorem, which
holds for every finite irreducible chain with weights $\prod P(u,\text{parent}(u))$, and is a probabilistic twin of the matrix-tree theorem of Tome~11, Appendix~A.
Wilson's algorithm (1996), with loop-erased walks, samples a uniform spanning tree in the expected time of a few hitting times instead of
a cover time.
\end{solution}

\begin{solution}{pb:terras}
\idea The parity of $T^k(m)$ is a coin flip that has not yet been decided by the first $k$ bits of $m$: each new bit of $m$ decides exactly
one new parity.

\writeup (a) We prove by induction on $k$ that $v_k(m)$ depends only on $m\bmod2^k$ and that $T^k(m+2^k)=T^k(m)+3^j$, where $j$ is the number of ones in $v_k(m)$. For
$k=1$: $m$ and $m+2$ have the same parity, and $T(m+2)=T(m)+1$ if $m$ is even, $T(m)+3$ if $m$ is odd. Suppose it holds for $k$, and let $m'=m+2^{k+1}$.
Applying the statement twice, $v_k(m')=v_k(m)$ and $T^k(m')=T^k(m)+2\cdot3^j$, which has the same parity as $T^k(m)$; so $v_{k+1}(m')=v_{k+1}(m)$, and
$T^{k+1}(m')=T^{k+1}(m)+3^{j'}$ where $j'$ counts the ones in $v_{k+1}(m)$ (apply $T$ to two numbers of the same parity differing by $2\cdot3^j$). This proves the
first claim. For the bijection, it suffices to show injectivity on residues modulo $2^k$. If $m'=m+2^{k-1}$, the induction shows that
$v_{k-1}(m')=v_{k-1}(m)$ and $T^{k-1}(m')=T^{k-1}(m)+3^{j}$, an odd difference: the last parities differ. So the two residues modulo $2^k$ that lift a residue
modulo $2^{k-1}$ have different vectors, and by induction all $2^k$ vectors are different.

(b) A run of $2^k$ consecutive integers contains each residue modulo $2^k$ once, so $v_k(m)$ is uniform on $\{0,1\}^k$: its coordinates are independent
fair bits.

(c) Induction: if $T^i(m)=\lambda_im+\rho_i$, the next step multiplies by $\frac12$, or by $\frac32$ and adds $\frac12$, so $T^{i+1}(m)=\lambda_{i+1}m+\rho_{i+1}$ with $\rho_{i+1}\ge0$
depending only on the parities so far.

(d) The vector $(1,1,1,1,1)$ corresponds to one residue modulo $32$; since $T(m)=\frac{3m+1}2$ keeps $m\equiv-1\pmod{2^r}$ in the form $-1\pmod{2^{r-1}}$, it is
$m\equiv-1$, and $31,47,71,107,161$ are all odd: $m=31$. The vector of $1$ is $(1,0,1,0,1)$ (the cycle $1,2,1,2,\dots$), so the next integer with this vector is
$1+32=33$, whose trajectory $33,50,25,38,19$ indeed alternates.

\tools{Tome~1, Chapter~1 (congruences, induction).}
\leads Terras proved (a) in 1976, and Everett found it independently in 1977; Lagarias's survey of 1985 and his annotated bibliography
are the standard guides to the $3x+1$ problem. The bijection makes the random model of Problem~\ref{pb:everett} an exact statement about
the first $k$ steps.
\end{solution}

\begin{solution}{pb:everett}
\idea If at least $k(1-\frac{\ln2}{\ln3})$ of the first $k$ parities are even, the $k$ steps have multiplied by less than $1$. By Problem~\ref{pb:terras} the parities are fair
coins on residue classes, and the law of large numbers does the rest.

\writeup (a) By Problem~\ref{pb:terras}(c), for $m$ with $v_k(m)=u$, $T^k(m)=\lambda m+\rho_u$ with $\lambda=\frac{3^j}{2^k}<1$. Then $T^k(m)<m$ as soon as $m>\frac{\rho_u}{1-\lambda}$, so $\sigma(m)\le k$.

(b) Let $U_k$ be the set of $u$ with $3^j<2^k$, and $r_u$ the residue modulo $2^k$ with $v_k(r_u)=u$ (Problem~\ref{pb:terras}(a)). By (a) there is $m_0$ such that every $m>m_0$ with
$m\bmod2^k\in\{r_u:u\in U_k\}$ has $\sigma(m)\le k$. The density of that union of residue classes is $\frac{|U_k|}{2^k}=P(3^{B_k}<2^k)=P(B_k<k\frac{\ln2}{\ln3})$. So the lower density of
$\{\sigma\le k\}$ is at least this.

(c) Since $\frac{\ln2}{\ln3}=0.6309\ldots>\frac12$, the weak law of large numbers (Tome~3, Theorem~4.2) gives $P(B_k<k\frac{\ln2}{\ln3})\to1$. The set $\{\sigma<\infty\}$ contains $\{\sigma\le k\}$
for every $k$, so its lower density is $1$, and its density exists and equals $1$.

(d) In the model, $\ln T^{i+1}(m)-\ln T^i(m)$ is $\ln\frac12$ or approximately $\ln\frac32$ (the $+\frac12$ is negligible for large values), with probability $\frac12$ each:
mean $\frac12\ln\frac34\approx-0.144$. A typical $m$ therefore needs about $\frac{\ln m}{\frac12\ln(4/3)}\approx6.95\ln m$ steps of $T$ to come down to $1$. For $m$ near $10^6$ this
predicts $96$ steps; the actual average over $10^6\le m<10^6+2000$ is $92.8$ (notebook). The prediction is a heuristic: the parities of later steps are
fair coins only in the sense of (a), for $k$ smaller than the number of digits of $m$.

\tools{Problem~\ref{pb:terras}; Tome~3, Theorem~4.2 (weak law of large numbers); natural density (Tome~1).}
\leads Tao proved in 2019 that almost all orbits (in logarithmic density) attain almost bounded values: for any function $f(m)\to\infty$, almost all $m$
have an iterate below $f(m)$. Whether every orbit reaches $1$ remains open; Tome~15, Appendix~F tells the story.
\end{solution}

\begin{solution}{pb:holte}
\idea Adding $m$ digits at a time is adding in base $b^m$, so the $m$-step chain is the one-step chain of a huge base, and in a huge base
the digits are uniform real numbers in disguise: the carry becomes $\lfloor U_1+\dots+U_n\rfloor$, whose law is Eulerian.

\writeup (a) $C_i$ is a function of $C_{i-1}$ and of the digits in position $i$, which are independent of $C_0,\dots,C_{i-1}$ (functions of earlier digits). So $(C_i)$ is a
Markov chain, with $P_b(i,j)=P(\lfloor(i+D_1+\dots+D_n)/b\rfloor=j)$ for independent uniform digits. The carry stays below $n$: if $C_{i-1}\le n-1$, then
$C_{i-1}+\sum D\le n-1+n(b-1)<nb$.

(b) Group positions $1,\dots,m$, $m+1,\dots,2m$, and so on: the $m$ digits of a block form one digit in base $b^m$, uniform on $\{0,\dots,b^m-1\}$, and the
carry out of a block, given the carry in, is computed by the same rule in base $B=b^m$. So the carries $C_m,C_{2m},\dots$ form the carry chain of
base $b^m$, and $P_b^m=P_{b^m}$.

(c) Couple: with $D_r=\lfloor BU_r\rfloor$, the digits are independent and uniform on $\{0,\dots,B-1\}$, and $\frac{i+\sum D_r}B$ differs from $S=\sum U_r$ by less than $\frac nB$. The two floors
can differ only if $S$ lies within $\frac nB$ of one of the integers $1,\dots,n-1$; as $S$ has a density at most $1$, this has probability at most $(n-1)\frac{2n}B$. So
the row $P_B(i,\cdot)$ and the law of $\lfloor S\rfloor$ are at total variation distance at most $\frac{2n^2}B$, and so is every entry.

(d) The law of $\lfloor S\rfloor$ is $\pi(j)=\frac{A(n,j)}{n!}$ by Problem~\ref{pb:irwinhall}. Let $\nu$ be a stationary law of $P_b$ (one exists, Tome~3, Theorem~8.8). Then
$\nu=\nu P_b^m=\nu P_{b^m}$, and every row of $P_{b^m}$ is within $\frac{2n^2}{b^m}$ of $\pi$; letting $m\to\infty$, $\nu=\pi$. The same coupling bound, for $B=b^m$, is the rate.

(e) For $n=3$, $\pi=(\frac16,\frac46,\frac16)$: the carry is $1$ in two thirds of the positions. In general $S$ and $n-S$ have the same law, so the mean of $\lfloor S\rfloor$
is $\frac{n-1}2$ (as $\lfloor n-S\rfloor=n-1-\lfloor S\rfloor$ almost surely).

\tools{Tome~3, Chapter~8 (Markov chains, total variation); Problems~\ref{pb:carries} and \ref{pb:irwinhall}.}
\leads Holte (1997) found the full eigenstructure: the eigenvalues are $1,b^{-1},\dots,b^{-(n-1)}$, and the left eigenvectors do not depend on $b$. Diaconis and
Fulman (2009) explained why: the carries chain has the same transition matrix as the number of descents after riffle shuffles of
Tome~3, Chapter~8, the base $b$ playing the role of the number of packets.
\end{solution}

\begin{solution}{pb:coprimewalk}
\idea Visibility is a sieve over primes. For each prime the walk equidistributes modulo $p$, except that $2$ is special: the parity of $X_t+Y_t$
is that of $t$. A uniform bound on the probability of divisibility by a large prime makes the infinite sieve converge.

\writeup (a) The steps $(\pm1,0)$ and $(0,\pm1)$ change $(U,V)$ by $(1,1)$, $(-1,-1)$, $(1,-1)$, $(-1,1)$, each with probability $\frac14$: independent fair signs.

(b) For odd $p$, $p\mid X$ and $p\mid Y$ if and only if $p\mid U$ and $p\mid V$, because $X=\frac{U+V}2$, $Y=\frac{U-V}2$ and $2$ is invertible modulo $p$. By independence the
probability is $P(p\mid U_t)^2$. The walk $U_t\bmod p$ on $\Z/p\Z$ with steps $\pm1$ is aperiodic for odd $p$ ($\{\pm1\}$ lies in no coset of the trivial subgroup, the only
proper one), so $P(p\mid U_t)\to\frac1p$ by Problem~\ref{pb:period}. For $p=2$: $X_t+Y_t=U_t\equiv t\pmod2$, so for odd $t$ one coordinate is odd. For even $t$, $U$ and $V$ are
even, and $X,Y$ are both even iff $U\equiv V\pmod4$. The walk $U_t\bmod4$ at even times moves on $\{0,2\}$ by steps $0,\pm2$ with probabilities
$\frac12,\frac14,\frac14$, so it converges to the uniform law on $\{0,2\}$, and so does $V$, independently: $P(U\equiv V\bmod4)\to\frac12$.

(c) $P(p\mid U_t)=\frac1p\sum_{k=0}^{p-1}E[e^{2\pi ikU_t/p}]=\frac1p\sum_{k=0}^{p-1}\cos^t\frac{2\pi k}p$. For $|y|\le\frac\pi2$, $\cos y=1-2\sin^2\frac y2\le1-\frac{2y^2}{\pi^2}\le e^{-2y^2/\pi^2}$, using $\sin\frac y2\ge\frac y\pi$; as $|\cos|$ has period
$\pi$, $|\cos x|\le e^{-2\delta(x)^2/\pi^2}$ where $\delta(x)$ is the distance from $x$ to $\pi\Z$. As $k$ runs over $\{1,\dots,p-1\}$, $\delta(\frac{2\pi k}p)$ takes each value $\frac{\pi j}p$, $1\le j\le\frac{p-1}2$, exactly
twice. So
\[
\sum_{k=0}^{p-1}\Bigl|\cos\frac{2\pi k}p\Bigr|^t\le1+2\sum_{j\ge1}e^{-2tj^2/p^2}\le1+2\int_0^\infty e^{-2tx^2/p^2}\dd x=1+p\sqrt{\frac\pi{2t}},
\]
and $P(p\mid U_t)\le\frac1p+\sqrt{\frac\pi{2t}}\le\frac1p+\frac2{\sqrt t}$.

(d) Let $E_p$ be the event that $p$ divides both coordinates; $\gcd(X_t,Y_t)=1$ means that no $E_p$ occurs and $(X_t,Y_t)\ne(0,0)$. Fix $P$ and let $Q=4\prod_{2<p\le P}p$. The pair
$(U_t\bmod Q,V_t\bmod Q)$ has independent coordinates, each of which is the $\pm1$ walk on $\Z/Q\Z$; this walk has period $2$ and at even (respectively odd) times converges to the uniform law on the
even (respectively odd) residues, by Problem~\ref{pb:period}. By the Chinese remainder theorem, along odd $t$ the events $E_p$, $p\le P$, become independent with
limits $P(E_2)=0$ and $P(E_p)=\frac1{p^2}$; along even $t$, $P(E_2)\to\frac12$ and $P(E_p)\to\frac1{p^2}$, again independently. So
\[
P\Bigl(\bigcap_{p\le P}E_p^c\Bigr)\to\alpha\prod_{2<p\le P}\Bigl(1-\frac1{p^2}\Bigr),\qquad\alpha=1\ (t\text{ odd}),\quad\alpha=\frac12\ (t\text{ even}).
\]
For the large primes: if $p>t$, $E_p$ forces $X_t=Y_t=0$, since $|X_t|,|Y_t|\le t$. For $P<p\le t$, by (c), $P(E_p)\le(\frac1p+\frac2{\sqrt t})^2\le\frac2{p^2}+\frac8t$. Hence
\[
0\le P\Bigl(\bigcap_{p\le P}E_p^c\Bigr)-P(\gcd(X_t,Y_t)=1)\le\sum_{p>P}\frac2{p^2}+\frac{8\pi(t)}t+P(X_t=Y_t=0),
\]
where $\pi(t)$ counts primes. As $t\to\infty$ the last two terms tend to $0$ (Chebyshev's bound gives $\pi(t)=O(\frac t{\ln t})$, Tome~1, Chapter~2), and then the first tends to $0$ as $P\to\infty$. Since $\prod_{p}(1-p^{-2})=\frac1{\zeta(2)}=\frac6{\pi^2}$ and the factor at $p=2$ is $\frac34$, $\prod_{p>2}(1-p^{-2})=\frac8{\pi^2}$. The limits are
$\frac8{\pi^2}\approx0.811$ for odd $t$ and $\frac4{\pi^2}\approx0.405$ for even $t$; the notebook finds $0.8108$ and $0.4055$ at $t=3001$ and $t=3000$.

\tools{Problem~\ref{pb:period}; Tome~1, Chapters~1--2 (Chinese remainder theorem, $\pi(t)=o(t)$, the Euler product for $\zeta(2)$); Tome~0 (Gaussian integral).}
\leads The average of the two limits is $\frac6{\pi^2}$, the density of visible points: the walk sees the plane as a uniform observer does, except for a
parity that it cannot forget. The same sieve works for any walk whose steps generate $\Z^2$, and the general phenomenon, an arithmetic
property of a random walk governed by the walk's behaviour modulo every integer, is the starting point of the theory of \emph{affine sieves}
of Bourgain, Gamburd and Sarnak.
\end{solution}

\begin{solution}{pb:range}
\idea A new point is discovered at time $s$ exactly when the walk, read backwards from time $s$, never comes back to its starting point.
So the expected range is a sum of non-return probabilities, which Tome~3 computed.

\writeup (a) $R_t=1+\sum_{s=1}^t\mathbf 1_{N_s}$, where $N_s$ is the event that $S_s\ne S_r$ for all $r<s$. Let $S'_q=S_s-S_{s-q}$ for $0\le q\le s$: it is again a simple walk (the steps in
reverse order), and $N_s$ says $S'_q\ne0$ for $1\le q\le s$. So $P(N_s)=P(S_1\ne0,\dots,S_s\ne0)$, and the term $s=0$ accounts for the $1$.

(b) Returns to $0$ only happen at even times, so the probability of no zero among times $1,\dots,2k+1$ equals that among times $1,\dots,2k$, which is
$u_{2k}$ by Tome~3, Theorem~7.3(a). Hence the terms $s=2k$ and $s=2k+1$ of (a) are both $u_{2k}$, and
\[
\mathbb E[R_{2m}]=2\sum_{k=0}^{m-1}u_{2k}+u_{2m}.
\]
Since $u_{2k+2}=u_{2k}\frac{2k+1}{2k+2}$, we have $(2m+1)u_{2m}=(2m+2)u_{2m+2}$, and the identity $\sum_{k=0}^mu_{2k}=(2m+1)u_{2m}$ follows by induction:
$(2m+1)u_{2m}+u_{2m+2}=(2m+3)u_{2m+2}$. So $\sum_{k=0}^{m-1}u_{2k}=2mu_{2m}$ and $\mathbb E[R_{2m}]=(4m+1)u_{2m}$. (Check: $m=1$ gives $\frac52$, the average of $2,3,3,2$.)

(c) $u_{2m}\sim\frac1{\sqrt{\pi m}}$ by Stirling's formula (Tome~0, Theorem~5.6), so
\[
\mathbb E[R_{2m}]\sim\frac{4m}{\sqrt{\pi m}}=\sqrt{\frac{16m}\pi}=\sqrt{\frac{8\cdot2m}\pi};
\]
odd times are squeezed between even ones since $R_t$ is non-decreasing.

(d) The same reversal gives
\[
\mathbb E[R_t]=\sum_{s=0}^tq_s,\qquad q_s=P(S_r\ne0\text{ for }1\le r\le s),
\]
and $q_s$ decreases to the probability of never returning. A Cesàro mean of a convergent sequence has the same limit.

\tools{Tome~3, Theorem~7.3 and Theorem~7.12; Stirling's formula (Tome~0, Theorem~5.6).}
\leads Dvoretzky and Erdős (1951) proved the law of large numbers $\frac{R_t}t\to1-F$ almost surely in $\Z^d$, $d\ge3$, and found the order $\frac t{\ln t}$ in the plane;
in $\Z^3$ the walk returns with probability about $0.34$, so it discovers about $0.66$ new points per step. The range is the simplest of the
\emph{occupation} questions that make random walks a laboratory for statistical physics.
\end{solution}

\begin{solution}{pb:cyclecover}
\idea The visited set grows as an interval, and each extension is a gambler's ruin. For the last vertex, only the two neighbours of the
candidate matter.

\writeup (a) The walker moves to adjacent vertices, so the visited set is connected: an arc. When it has $k<n$ vertices, the walker is at one of its
ends (it has just added it) or, at time $0$, at the unique vertex. Label the arc $1,\dots,k$ with the walker at $1$, and the two unvisited
neighbouring vertices $0$ and $k+1$ (distinct since $k<n-1$; if $k=n-1$ they coincide, which changes nothing below). A new vertex is added when the
walker leaves $\{1,\dots,k\}$, which takes $1\cdot(k+1-1)=k$ steps on average (Tome~3, Theorem~7.9).

(b) $\sum_{k=1}^{n-1}k=\frac{n(n-1)}2$.

(c) Fix $v\ne0$. Before it visits $v$, the walker must visit $v-1$ or $v+1$ (it starts at $0\ne v$, and $0$ may be one of them); let $T$ be the first time it is at
one of them, say at $v-1$ (the other case is symmetric). At time $T$, $v$ is unvisited. Now $v$ is the last vertex visited if and only if, from $v-1$,
the walker reaches $v+1$ before $v$, going the long way round. By the gambler's ruin on the arc from $v$ to $v$ through $v-1,\dots,v+1$, of length $n$, this has
probability $\frac1{n-1}$, by the strong Markov property at $T$.

(d) On $K_n$ the walk from $0$ is invariant under permutations of the other vertices, so the last one is uniform. On the star with centre $c$ and leaves
$\ell_1$, $\ell_2$, $\ell_3$, started at $\ell_1$, the centre is visited at time $1$, so the last vertex is $\ell_2$ or $\ell_3$, never $c$.

\tools{Tome~3, Theorem~7.9 (gambler's ruin) and the strong Markov property (Chapter~8).}
\leads Lovász and Winkler showed in 1993 that the cycle and the complete graph are the only graphs on which the last vertex is uniform from every
start. Cover times of large graphs are a deep subject: on the $n\times n$ torus the cover time is about $\frac4\pi n^2(\ln n)^2$ (Dembo, Peres, Rosen and
Zeitouni, 2004).
\end{solution}

\begin{solution}{pb:threegamblers}
\idea Polynomials in the fortunes whose expected change is constant turn expected durations into algebra. The product $XYZ$ loses on
average one third of the total per round, and the sum of pairwise products loses exactly one per round.

\writeup (a) Suppose all three are in the game. If the pair $(X,Y)$ plays, the change is $(X+\varepsilon)(Y-\varepsilon)Z-XYZ=(\varepsilon(Y-X)-\varepsilon^2)Z$ with $\varepsilon=\pm1$ equally likely,
of conditional mean $-Z$. Averaging over the three pairs, $\mathbb E[\Delta(XYZ)\mid\mathcal F_t]=-\frac{X+Y+Z}3=-\frac s3$. So $M_t=X_tY_tZ_t+\frac s3t$, stopped at the first ruin
$\tau$, is a martingale. The state space is finite and from every state a ruin can be reached within $s$ rounds, so $\mathbb E[\tau]<\infty$ (Tome~3, Lemma~7.8), and the
increments are bounded; optional stopping (Tome~3, Theorem~10.12) gives $abc=\mathbb E[M_\tau]=0+\frac s3\mathbb E[\tau]$.

(b) With three players, the pair $(X,Y)$ changes $XY+YZ+ZX$ by $\varepsilon(Y-X)-\varepsilon^2+\varepsilon Z-\varepsilon Z$, of mean $-1$, and similarly for the other pairs. With two players,
say $Z=0$, only $XY$ changes, by $\varepsilon(Y-X)-1$, of mean $-1$. So $Q_t+t$ is a martingale until the end $\sigma$, where $Q=XY+YZ+ZX$ is $0$. As in (a), $\mathbb E[\sigma]<\infty$ and
$\mathbb E[\sigma]=Q_0=ab+bc+ca$.

(c) $X_t$ is a bounded martingale (each round is fair for each player, whether or not the player plays), so $a=\mathbb E[X_\sigma]=s\,P(\text{first player wins all})$.

(d) With $a=b=c=4$: the game lasts $48$ rounds on average, and the first player leaves after $16$ rounds on average.

\tools{Tome~3, Theorem~10.12 (optional stopping) and Lemma~7.8.}
\leads For two players the same martingale $XY+t$ is the gambler's ruin duration $ab$ of Tome~3, Theorem~7.9. The probability that a given player is
the first to be ruined is a much harder quantity: it is a harmonic measure for a random walk in a triangle. The analogue with $n$ players is the expected time to absorption of a random walk in a simplex.
\end{solution}

\begin{solution}{pb:moran}
\idea Each step is fair, so the type frequency is a martingale; the only randomness is in the speed, which is slowest when one type is
rare. Counting moves state by state reduces the time to a Green function of the simple walk.

\writeup (a) From $x$, the walk goes up if an $A$ reproduces and a $B$ dies, probability $\frac{x(N-x)}{N^2}$, and down with the same probability. So $X$ is a bounded
martingale, absorbed at $0$ or $N$ in finite expected time (Tome~3, Lemma~7.8), and $i=\mathbb E[X_\tau]=N\,P(X_\tau=N)$.

(b) With $\delta=X_{t+1}-X_t$: $X_{t+1}(N-X_{t+1})=X_t(N-X_t)+\delta(N-2X_t)-\delta^2$, and $\mathbb E[\delta\mid X_t]=0$, $\mathbb E[\delta^2\mid X_t]=\frac{2X_t(N-X_t)}{N^2}$.

(c) Let $q_x=\frac{2x(N-x)}{N^2}$, the probability of moving from $x$. The sequence of positions after each move is the simple random walk on $\{0,\dots,N\}$ absorbed at the ends,
and each stay at $x$ before a move lasts on average $\frac1{q_x}$ steps (geometric). Let $G(i,x)$ be the expected number of moves from $x$ made by the simple walk
started at $i$, that is its expected number of visits to $x$. For $x\ge i$ the walk reaches $x$ with probability $\frac ix$ (gambler's ruin). From $x$ it escapes to $\{0,N\}$
before returning to $x$ with probability $\frac12\bigl(\frac1x+\frac1{N-x}\bigr)=\frac N{2x(N-x)}$, so the number of visits, once $x$ is reached, is geometric with mean $\frac{2x(N-x)}N$. Multiplying, $G(i,x)=\frac ix\cdot\frac{2x(N-x)}N=\frac{2i(N-x)}N$ for $x\ge i$, and symmetrically $G(i,x)=\frac{2x(N-i)}N$ for $x\le i$. Hence
\[
\begin{aligned}t_i=\sum_{x=1}^{N-1}\frac{G(i,x)}{q_x}&=\sum_{x\ge i}\frac{2i(N-x)}N\cdot\frac{N^2}{2x(N-x)}+\sum_{x<i}\frac{2x(N-i)}N\cdot\frac{N^2}{2x(N-x)}\\&=Ni\sum_{x=i}^{N-1}\frac1x+N(N-i)\sum_{x=1}^{i-1}\frac1{N-x},\end{aligned}
\]
and the last sum is $\sum_{y=N-i+1}^{N-1}\frac1y$.

(d) For $i=1$ the second sum is empty and $t_1=NH_{N-1}\sim N\ln N$. For $i=\lfloor pN\rfloor$, $\sum_{x=i}^{N-1}\frac1x=\ln\frac Ni+O(\frac1i)=\ln\frac1p+o(1)$ and $\sum_{y=N-i+1}^{N-1}\frac1y=\ln\frac1{1-p}+o(1)$, so
$t_i=N^2\bigl(p\ln\frac1p+(1-p)\ln\frac1{1-p}+o(1)\bigr)$.

\tools{Tome~3, Theorems~7.9 and 10.12, Lemma~7.8; harmonic sums (Tome~0).}
\leads Moran introduced the model in 1958 as a tractable cousin of the Wright--Fisher model, in which the whole generation is resampled at
once; the diffusion approximation of Wright and Kimura gives the same entropy-like formula for that model, with time counted in
generations. The decay of heterozygosity in (b) is the population geneticist's measure of genetic drift. Kingman's coalescent (1982) runs the same
models backwards in time, following the ancestry of a sample.
\end{solution}

\begin{solution}{pb:bernoullilaplace}
\idea The chain maps polynomials of degree $k$ to polynomials of degree $k$, so its matrix is triangular in the basis of powers, and the
eigenvalues are read on the diagonal.

\writeup (a) From $x$: the number of red balls in the first urn drops if a red is drawn there ($\frac xn$) and a white in the second (the second urn
has $n-x$ red and $x$ white): $P(x,x-1)=\frac{x^2}{n^2}$. It rises with probability $P(x,x+1)=\frac{(n-x)^2}{n^2}$, and stays otherwise. Detailed balance:
$\pi(x)P(x,x+1)=\binom nx^2\frac{(n-x)^2}{n^2}\binom{2n}n^{-1}$ and $\pi(x+1)P(x+1,x)=\binom n{x+1}^2\frac{(x+1)^2}{n^2}\binom{2n}n^{-1}$, equal since $\binom nx(n-x)=\binom n{x+1}(x+1)$; and $\sum_x\binom nx^2=\binom{2n}n$.

(b) With $\delta=X_1-x\in\{-1,0,1\}$: $\mathbb E[\delta]=\frac{(n-x)^2-x^2}{n^2}=\frac{n-2x}n$ and $\mathbb E[\delta^2]=\frac{x^2+(n-x)^2}{n^2}$. Then $\mathbb E[(x+\delta)^k]=\sum_r\binom krx^{k-r}\mathbb E[\delta^r]$, and $\mathbb E[\delta^r]$ is
$\mathbb E[\delta]$ or $\mathbb E[\delta^2]$, a polynomial in $x$ of degree $1$ or $2$. So the terms $r\ge3$ have degree at most $k-1$, and the coefficient of $x^k$ in $P(x^k)$ is
$1+k\cdot\bigl(-\frac2n\bigr)+\binom k2\frac2{n^2}=1-\frac{2k}n+\frac{k(k-1)}{n^2}=1-\frac{k(2n-k+1)}{n^2}$. Functions on $\{0,\dots,n\}$ are exactly the polynomials of degree at most $n$ (Lagrange interpolation),
and in the basis $1,x,\dots,x^n$ the operator $P$ is triangular with diagonal $\lambda_0,\dots,\lambda_n$. These are its eigenvalues.

(c) $\mathbb E[X_{t+1}-\frac n2\mid X_t]=X_t-\frac n2+\frac{n-2X_t}n=(1-\frac2n)(X_t-\frac n2)$, so $\mathbb E[X_t]-\frac n2=\frac n2(1-\frac2n)^t$. For $0\le y\le\frac12$, $\ln(1-y)\ge-y-y^2$, so for $n\ge4$ and $t\le\frac n4\ln n$,
$t\ln(1-\frac2n)\ge-\frac n4\ln n\,(\frac2n+\frac4{n^2})=-\frac12\ln n-\frac{\ln n}n$, and $\frac n2(1-\frac2n)^t\ge\frac{\sqrt n}2e^{-\ln n/n}\ge\frac{\sqrt n}2e^{-1/e}>\frac{\sqrt n}3$. Under $\pi$, $X$ is hypergeometric (draw $n$
balls from $2n$, half red), with variance $\frac{n^2}{4(2n-1)}<\frac n4$.

\tools{Tome~3, Proposition~8.14 (detailed balance); polynomial interpolation (Tome~0).}
\leads Daniel Bernoulli posed the model in 1769 and Laplace solved its diffusion limit in his \emph{Théorie analytique des probabilités} (1812). Diaconis
and Shahshahani (1987) used the eigenvalues above (with their eigenvectors, the dual Hahn polynomials) to show that the chain mixes after
$\frac n4\ln n$ steps and not before, a cutoff; (c) is the lower half of their theorem.
\end{solution}

\begin{solution}{pb:kacring}
\idea Each ball's colour is a product of independent signs along the window of edges it has crossed. Products over overlapping windows
are correlated, over disjoint windows independent, and there are only $O(t)$ overlapping windows per ball: the law of large numbers wins
for every fixed $t$, while the recurrence at time $2n$ moves away as $n$ grows.

\writeup (a) Backwards, every ball moves from $j+1$ to $j$ and changes colour if edge $j$ is marked; this undoes a tick. After $n$ ticks every ball has
crossed each edge once, so it has changed colour $M$ times and $\Delta_n=(-1)^M\Delta_0$; after $2n$ ticks it is on its site again, having changed colour $2M$ times.

(b) The ball starting at $i$ crosses edges $i,i+1,\dots,i+t-1$ during the first $t$ ticks, hence its colour. For $t\le n$ these edges are distinct, so
$\mathbb E[\varepsilon_i\cdots\varepsilon_{i+t-1}]=(1-2\mu)^t$, and $\mathbb E[\Delta_t]=(1-2\mu)^t\sum_ic_i$.

(c) Let $W_i=\{i,\dots,i+t-1\}$ and $\Pi_i=\prod_{j\in W_i}\varepsilon_j$. If $W_i\cap W_{i'}=\emptyset$, $\Pi_i$ and $\Pi_{i'}$ are independent and uncorrelated. For $t\le\frac{n+1}2$, $W_{i'}$ meets $W_i$ only when
the cyclic distance from $i$ to $i'$ is less than $t$, which happens for $2t-1$ values of $i'$. As $|\Pi_i|=1$, every covariance is at most $1$ in absolute value. So
$\operatorname{Var}(\Delta_t)=\sum_{i,i'}c_ic_{i'}\operatorname{Cov}(\Pi_i,\Pi_{i'})\le n(2t-1)$. By Chebyshev, $\frac{\Delta_t}n-(1-2\mu)^t\frac{\Delta_0}n\to0$ in probability; with $\frac{\Delta_0}n=2\beta-1$, the proportion of black
balls $\frac12(1+\frac{\Delta_t}n)$ tends to $\frac12+(1-2\mu)^t(\beta-\frac12)$.

(d) Part (c) is a statement for each fixed $t$ as $n\to\infty$ (and remains useful for $t$ small compared with $n$), while the returns of (a) happen at
times $n$ and $2n$, which run away to infinity. The variance bound $n(2t-1)$ is useless for $t$ of order $n$; at $t=n$ all windows are the whole
circle, all products equal $(-1)^M$, and the fluctuations are as large as possible. This is Boltzmann's answer to Zermelo: the
recurrence time of a gas is astronomically long. Loschmidt's reversal is answered by (c) as well: running the ring backwards from a
typical state at time $t$ is the dynamics of the ring with mirrored marks, but from a very special state, the only one that leads
back to all black. In the think box: with $3$ marked edges out of $6$, after $6$ ticks every ball has changed colour $3$ times, so all balls are
white; after $12$ ticks they are all black again. With $1000$ sites and $\mu=\frac1{10}$, after $10$ ticks the proportion of black balls is close to
$\frac12(1+0.8^{10})\approx0.554$, and after $1000$ ticks it is exactly $0$ or $1$, according to the parity of $M$.

\tools{Tome~3, Chapter~4 (variance, Chebyshev's inequality); Tome~2, Chapter~1 (the rotation group of the ring).}
\leads Gottwald and Oliver's article \emph{Boltzmann's dilemma} (2009) works out the Kac ring in detail and uses it as an introduction to
statistical mechanics. Kac's own account is in his \emph{Probability and Related Topics in Physical Sciences} (1959).
\end{solution}

\begin{solution}{pb:gilbreath}
\idea Both descriptions say that the cards seen from the top always form a block of consecutive numbers: a Gilbreath shuffle eats the
deck from the middle outwards.

\writeup (a) Call an arrangement \emph{interval-like} if for every $j$ its top $j$ cards form a block of consecutive integers. After dealing $k$ cards, the
pile reads $k,k-1,\dots,1$ from the top and the rest $k+1,\dots,n$; in any riffle the top $j$ cards are the top $r$ cards of the pile and the top $j-r$
of the rest, $\{k-r+1,\dots,k+j-r\}$, a block. Conversely, in an interval-like arrangement with top card $x$, each further card extends the current
block by one, just below it or just above it. So the cards smaller than $x$ appear in decreasing order and the cards larger than $x$ in increasing
order, and the arrangement is a riffle of the pile $x,x-1,\dots,1$ (deal $k=x$ cards) with the rest $x+1,\dots,n$. The Gilbreath arrangements are exactly the
interval-like ones. Reading from the bottom, an interval-like arrangement chooses at each step the smallest or the largest of the cards not yet
placed (the top $j$ cards form a block, so the card below them is just below or just above it), and the top card is forced: there are $2^{n-1}$ of
them.

Residues: a block of $j$ consecutive integers has distinct residues modulo $j$, so interval-like arrangements qualify. Conversely, suppose the top $j$ cards
have distinct residues modulo $j$ for every $j$. For $j=n-1$ the top $n-1$ cards are all of $\{1,\dots,n\}$ but one card $y$, with distinct residues modulo
$n-1$; as $1\equiv n\pmod{n-1}$, they cannot contain both $1$ and $n$, so $y\in\{1,n\}$. Remove the bottom card $y$; the remaining cards are $\{2,\dots,n\}$ or $\{1,\dots,n-1\}$,
and after subtracting $1$ in the first case (which preserves distinctness of residues) they again satisfy the condition. By induction they
are interval-like, and adding $y$ at the bottom keeps the arrangement interval-like.

(b) The suit of card $x$ depends only on $x\bmod4$. Among the top $4m$ cards the residues modulo $4m$ are distinct, so they are all the residues
modulo $4m$, and each residue modulo $4$ appears exactly $m$ times. The same holds for the top $4(m+1)$ cards, so the block of cards $4m+1,\dots,4m+4$
contains each residue modulo $4$, that is each suit, once. The same argument with $2$ gives the colours.

(c) Fix a Gilbreath arrangement with top card $x$. If $x$ came from the pile, $x=k$; if from the rest, $x=k+1$. In either case, given $k$, the
arrangement determines the interleaving: a card $\le k$ comes from the pile and a card $>k$ from the rest. Both values $k=x$ and $k=x-1$ work (when
they lie in $\{0,\dots,n\}$; for $x=1$, $k=0$ or $1$; for $x=n$, $k=n-1$ or $n$): reading the arrangement from the top, the cards $\le k$ appear in decreasing
order and the cards $>k$ in increasing order, because each new card extends the block below or above. So each arrangement arises from
exactly two pairs ($k$, interleaving) out of the $\sum_k\binom nk=2^n$ pairs, and if $k$ is binomial and the interleaving uniform, each pair has probability
$\binom nk2^{-n}\binom nk^{-1}=2^{-n}$: the arrangement is uniform, with probability $2^{1-n}$ each.

\tools{Tome~2, Chapter~4 (permutations); Tome~0 (counting, induction); Tome~3, Chapter~8 (riffle shuffles).}
\leads Norman Gilbreath described the principle in 1958 in the magicians' journal \emph{The Linking Ring}, and Martin Gardner popularised it
in his \emph{Scientific American} column. Diaconis and Graham's \emph{Magical Mathematics} (2012) gives the equivalence of (a), several further
characterisations, and many tricks built on it.
\end{solution}

\section*{Map of the chapter}

\begin{center}
\small
\adjustbox{max width=\linewidth}{\begin{tabular}{lll}
Problem & Topic & Tomes used\\\hline
\ref{pb:returnbound} & return probabilities on a group & 2, 3\\
\ref{pb:period} & periodicity and cosets & 1, 2, 3\\
\ref{pb:rollingdie} & a die rolling on a chessboard & 2, 3\\
\ref{pb:dodecaturns} & random turns of an icosahedron & 2, 3\\
\ref{pb:gapdiam} & spectral gap and diameter & 0, 2, 3\\
\ref{pb:treewalk} & the walk on the free group & 0, 3\\
\ref{pb:commute} & commute times and resistance & 3\\
\ref{pb:foster} & Foster's theorem, symmetric graphs & 0, 2, 3\\
\ref{pb:antipode} & hitting the opposite vertex & 2, 3\\
\ref{pb:rayleigh} & Thomson, Rayleigh, Nash-Williams & 3, 4\\
\ref{pb:aldousbroder} & uniform spanning trees & 3\\
\ref{pb:terras} & Collatz parity vectors & 1\\
\ref{pb:everett} & almost every orbit goes down & 1, 3\\
\ref{pb:holte} & the carries chain, Eulerian law & 3, 4\\
\ref{pb:coprimewalk} & visible points seen by a walk & 1, 2, 3\\
\ref{pb:range} & the range of a walk & 0, 3\\
\ref{pb:cyclecover} & cover time and last vertex of a cycle & 3\\
\ref{pb:threegamblers} & three gamblers & 3\\
\ref{pb:moran} & the Moran model & 0, 3, 4\\
\ref{pb:bernoullilaplace} & Bernoulli's two urns & 0, 3\\
\ref{pb:kacring} & Kac's ring & 2, 3\\
\ref{pb:gilbreath} & Gilbreath shuffles & 0, 2, 3\\
\end{tabular}}
\end{center}

\section*{Sources and further reading}

\begin{itemize}
\item H.~Poincaré, \emph{Science et méthode}, Flammarion, Paris, 1908; English translation by F.~Maitland, \emph{Science and Method}, Nelson, London,
1914 (the epigraph, from the chapter ``Chance'').
\item L.~Boltzmann, Weitere Studien über das Wärmegleichgewicht unter Gasmolekülen, \emph{Sitzungsberichte der Akademie der Wissenschaften
in Wien} 66 (1872), 275--370; J.~Loschmidt, Über den Zustand des Wärmegleichgewichtes eines Systems von Körpern mit Rücksicht auf die
Schwerkraft, \emph{ibid.} 73 (1876), 128--142; H.~Poincaré, Sur le problème des trois corps et les équations de la dynamique, \emph{Acta
Mathematica} 13 (1890), 1--270; E.~Zermelo, Ueber einen Satz der Dynamik und die mechanische Wärmetheorie, \emph{Annalen der Physik} 57
(1896), 485--494.
\item M.~Kac, Some remarks on the use of probability in classical statistical mechanics, \emph{Bulletins de l'Académie Royale de Belgique} 42
(1956), 356--361; M.~Kac, \emph{Probability and Related Topics in Physical Sciences}, Interscience, 1959; G.~A.~Gottwald and M.~Oliver,
Boltzmann's dilemma: an introduction to statistical mechanics via the Kac ring, \emph{SIAM Review} 51 (2009), 613--635.
\item P.~Diaconis, \emph{Group Representations in Probability and Statistics}, IMS Lecture Notes 11, 1988; P.~Diaconis and M.~Shahshahani,
Generating a random permutation with random transpositions, \emph{Zeitschrift für Wahrscheinlichkeitstheorie} 57 (1981), 159--179.
\item C.~Hooley, On Artin's conjecture, \emph{Journal für die reine und angewandte Mathematik} 225 (1967), 209--220.
\item P.~Diaconis and L.~Saloff-Coste, Comparison theorems for reversible Markov chains, \emph{Annals of Applied Probability} 3 (1993),
696--730; D.~B.~Wilson, Mixing times of lozenge tiling and card shuffling Markov chains, \emph{Annals of Applied Probability} 14 (2004),
274--325.
\item H.~Kesten, Symmetric random walks on groups, \emph{Transactions of the AMS} 92 (1959), 336--354.
\item A.~K.~Chandra, P.~Raghavan, W.~L.~Ruzzo, R.~Smolensky and P.~Tiwari, The electrical resistance of a graph captures its commute and
cover times, \emph{Proceedings of the 21st ACM Symposium on Theory of Computing} (1989), 574--586; R.~Aleliunas, R.~M.~Karp, R.~J.~Lipton,
L.~Lovász and C.~Rackoff, Random walks, universal traversal sequences, and the complexity of maze problems, \emph{Proceedings of the 20th
IEEE Symposium on Foundations of Computer Science} (1979), 218--223; O.~Reingold, Undirected connectivity in log-space, \emph{Journal of the ACM}
55 (2008), article 17.
\item R.~M.~Foster, The average impedance of an electrical network, in \emph{Contributions to Applied Mechanics} (Reissner Anniversary
Volume), Edwards, Ann Arbor, 1949; A.~E.~Brouwer, A.~M.~Cohen and A.~Neumaier, \emph{Distance-Regular Graphs}, Springer, 1989.
\item C.~St.~J.~A.~Nash-Williams, Random walk and electric currents in networks, \emph{Proceedings of the Cambridge Philosophical Society} 55
(1959), 181--194; P.~G.~Doyle and J.~L.~Snell, \emph{Random Walks and Electric Networks}, MAA, 1984 (free online); R.~Lyons and Y.~Peres,
\emph{Probability on Trees and Networks}, Cambridge University Press, 2016.
\item A.~Broder, Generating random spanning trees, \emph{Proceedings of the 30th IEEE Symposium on Foundations of Computer Science} (1989),
442--447; D.~J.~Aldous, The random walk construction of uniform spanning trees and uniform labelled trees, \emph{SIAM Journal on Discrete
Mathematics} 3 (1990), 450--465; D.~B.~Wilson, Generating random spanning trees more quickly than the cover time, \emph{Proceedings of the 28th
ACM Symposium on Theory of Computing} (1996), 296--303.
\item R.~Terras, A stopping time problem on the positive integers, \emph{Acta Arithmetica} 30 (1976), 241--252; C.~J.~Everett, Iteration of the
number-theoretic function $f(2n)=n$, $f(2n+1)=3n+2$, \emph{Advances in Mathematics} 25 (1977), 42--45; J.~C.~Lagarias, The $3x+1$ problem and its
generalizations, \emph{American Mathematical Monthly} 92 (1985), 3--23; T.~Tao, Almost all orbits of the Collatz map attain almost bounded
values, \emph{Forum of Mathematics, Pi} 10 (2022), e12.
\item J.~M.~Holte, Carries, combinatorics, and an amazing matrix, \emph{American Mathematical Monthly} 104 (1997), 138--149; P.~Diaconis and
J.~Fulman, Carries, shuffling, and symmetric functions, \emph{Advances in Applied Mathematics} 43 (2009), 176--196.
\item A.~Dvoretzky and P.~Erdős, Some problems on random walk in space, \emph{Proceedings of the Second Berkeley Symposium on Mathematical
Statistics and Probability} (1951), 353--367.
\item L.~Lovász and P.~Winkler, A note on the last new vertex visited by a random walk, \emph{Journal of Graph Theory} 17 (1993), 593--596; A.~Dembo,
Y.~Peres, J.~Rosen and O.~Zeitouni, Cover times for Brownian motion and random walks in two dimensions, \emph{Annals of Mathematics} 160
(2004), 433--464.
\item P.~A.~P.~Moran, Random processes in genetics, \emph{Proceedings of the Cambridge Philosophical Society} 54 (1958), 60--71; J.~F.~C.~Kingman,
The coalescent, \emph{Stochastic Processes and their Applications} 13 (1982), 235--248.
\item P.~Diaconis and M.~Shahshahani, Time to reach stationarity in the Bernoulli--Laplace diffusion model, \emph{SIAM Journal on
Mathematical Analysis} 18 (1987), 208--218.
\item N.~Gilbreath, Magnetic colors, \emph{The Linking Ring} 38 (1958); P.~Diaconis and R.~Graham, \emph{Magical Mathematics}, Princeton
University Press, 2012.
\end{itemize}
